% =====================================================================
%  R. Nielsen, ON LIFE AND HOPE.  Six Lectures.  Copenhagen: Gyldendal (F. Hegel), 1875.  103 pp.
%  (Om Liv og Haab.  Sex Forelæsninger.)
%  ENGLISH TRANSLATION, from transcription.tex (Danish).
%  GENERATED by ebook-method/liv-og-haab/tr_build.py --write from this template and the six translators' parts
%  in ebook-method/liv-og-haab/tr/ (TRANSLATOR-BRIEF.md; terms: GLOSSARY.md, the title and two key terms
%  approved by Hans 2026-10-12).  Conventions: TRANSLATION-PLAYBOOK.md.  Page numbers of the 1875 print in the margin.
% =====================================================================
\documentclass[12pt,a4paper,openany]{book}
\usepackage[utf8]{inputenc}
\usepackage[T1]{fontenc}
\usepackage{libertinus}
\usepackage{libertinust1math}
\usepackage[english]{babel}
\usepackage[protrusion=true,expansion=true]{microtype}
\usepackage{setspace}
\usepackage[margin=1.3in]{geometry}
\usepackage{fancyhdr}
\usepackage{hyperref}
\hypersetup{pdftitle={On Life and Hope}, pdfauthor={Rasmus Nielsen}, hidelinks}
\setlength{\headheight}{15pt}
\pagestyle{fancy}
\fancyhf{}
\fancyhead[LE,RO]{\thepage}
\fancyhead[C]{\textit{R. Nielsen: On Life and Hope (1875)}}
\setstretch{1.3}
\usepackage{marginnote}
\renewcommand{\thefootnote}{*)}
\newcommand{\opage}[1]{\marginnote{\footnotesize\textit{[#1]}}}
% a lecture opens a new page: heading, a short rule
\newcommand{\lecture}[1]{\clearpage\phantomsection\addcontentsline{toc}{chapter}{#1}\begin{center}{\large #1}\\[0.4em]\rule{3em}{0.4pt}\end{center}\medskip\noindent}

\begin{document}

% PDF 11, the title page
\begin{titlepage}
\begin{center}
\vspace*{2em}
{\small ON}\\[1em]
{\Huge LIFE AND HOPE.}\\[2.5em]
SIX LECTURES\\[1em]
{\small BY}\\[1em]
{\large R. NIELSEN.}\\[5em]
{\large COPENHAGEN.}\\[0.5em]
PUBLISHED BY THE GYLDENDAL BOOKSHOP (F. HEGEL).\\[0.5em]
{\small PRINTED BY J. JØRGENSEN \& CO.}\\[0.5em]
{\small 1875.}\\[3em]
{\small Translated from the Danish}
\end{center}
\end{titlepage}

% the print has no table of contents; this one is editorial
\renewcommand{\contentsname}{Contents {\normalsize\textit{(editorial; not in the print)}}}
\tableofcontents
\clearpage

% --- p. 1 ---
\setcounter{footnote}{0}%
\lecture{First Lecture.}\opage{1}Life, and particularly human life, is laid out on peculiar
conditions. We live in the present, and the present must be important to us, for it is indeed what
is real in the fullest measure. But if we now fix our gaze on this real, present thing, it shrinks
and becomes so short: it transforms itself into nothing but fleeting, passing moments. What in this
now is something present, is in the next something past, and what is still something to come, is in
the next now something present. If the whole content of our life were to be restricted to what we
must, in the strictest sense, call the present, how vain, empty and poor it would be. The present is
indeed really only a boundary between past and future; if we follow the movement of the boundary
point as a line, then the line of reality is narrow, narrower still than the one on which the
tightrope dancer moves.

How does it come, then, that the reality of life can not only spread itself out for our view in
space, but stretch itself in time, so that not merely minutes and hours, but days, months, years can
seem to gather into one and come under the present? It comes, of course, from this: that we are not
merely sunk in the temporal
% --- p. 2 ---
\setcounter{footnote}{0}%
\opage{2}now, but also raised above it: we look forward and we look back. We have what is coming in
expectation, what is gone in memory. Memory and expectation reach each other a hand and bear the
weight of reality on strong arms. Expectation of an evil is fear, expectation of a good is hope;
thus life is divided between fear and hope. Expectation aims at the possible, but memory explains
the real; for just as the plant grows from the seed, and the seed again from the plant, so the real
is always born of a possibility, and new possibilities lie hidden like seeds in the soil of reality.

How expectation alternates between fear and hope depends not only on something visible to the eye
that awakens dread, as in a shipwreck, when the cry of distress is heard and the lifeboats are at
work; but it depends also, and chiefly, on a deeper understanding of what reality is in its ground,
what possibilities it bears in its womb. Such an understanding depends in turn on experience, and
what has been experienced is preserved in memory: from memory we learn to fear, and from memory we
learn to hope.

When the year changes, there is a midnight hour when it can be suspenseful enough to be alone with
oneself. A curtain of the future rolls up, and what is it we see? A sea of billowing possibilities!
Otherwise the stage is empty, what moves across it is only shadows: the real that is expected has
not yet come. And yet the empty stage is built on real ground, it does not hover in the air: its
real ground is all that is gone, which memory bears. What will appear? The wish tempts, fear warns,
probability reckons with many figures, memory weaves a
% --- p. 3 ---
\setcounter{footnote}{0}%
\opage{3}riddling web, but hope? If hope will not appear, then how dead and desolate is the whole,
the breathing heavy, the stone on the heart not to be borne. But it grows light: Hope comes as the
herald of life, a dawning in the east, and yet behind the veil of memory like the moon in the clear
night. It is Hope with which the play of the future begins, it is Hope that is greeted with secret
jubilation. Hope is the beautiful actress: what a radiance in her glance! her words are music, she
dances on the waves of time, she casts --- see, these flowers of light on the sea of possibilities
across the dim stage peopled by shadows.

Although memory belongs to the past and hope to the future, it is nonetheless unmistakable that they
are two streams from one source. Life is the common source from which they well forth and to which
they unceasingly return. If the spring of hope dries up, then that of memory dries up too, and
conversely.

To judge by the outward, one would probably suppose that hope and memory were divided to such a
degree that they must even have each its own age of life. Youth, they say, is the time of hope, old
age the time of memory. The young person has little to remember, but much to hope; the old person
has little to hope, but much to remember. So one praises the spring of youth with its budding hopes
and consoles old age with an aftermath of memory. A beguiling eulogy, a sorry consolation. The hope
of youth that lacks memory, because it lacks the seal of (practical) experience upon inner
experiences, what is it, really? A fragrance, a sound, a run of notes on the stringed instrument of
life, a morning dream that vanishes when reality itself begins its day's work. And the memory
% --- p. 4 ---
\setcounter{footnote}{0}%
\opage{4}that is to console an old age without hope, how pale and faint it is! There is life in
memory as long as hope lives, but when hope dies, memory changes color. It sickens, wastes away, a
remembrance of loss, a ghost from a vanished time, a restless spirit whose staring gaze and
melancholy looks betray that there is not even peace in the graves of oblivion.

Life is more than merely a gift of nature; it is, spiritually understood, a great task. We first
live when we learn to live. Hope and memory do not belong to the things that present themselves of
their own accord without further ado; they are activities of the soul, inner states, which need care
and tending. Here there is something to learn. We learn to hope as we learn to fear, and we learn to
fear and to hope as we learn to remember.

Just as a good householder carefully goes through his account books and watches closely that no
items of debt are overlooked and nothing outstanding is forgotten, so must the running account of
our life between hope and memory continually be settled: outstanding sums collected, items of debt
liquidated, and memory must exactly write down and keep every receipt. To recollect is not the same
as to remember. There are people who remember well but recollect poorly, and there are those who
remember poorly but recollect well. Memory preserves the real simply as it is: we remember odds and
ends, the accidental with the essential; but recollection undertakes a purification of what has been
experienced and appropriates only the essential. Recollection fixes the capital, the running
expenses are defrayed by memory: in order to save up a capital of life, we must learn to recollect.
So too with hope. To hope is not the same as to
% --- p. 5 ---
\setcounter{footnote}{0}%
\opage{5}abound in glad expectation or confidently to wish. There are lively people, rosy-cheeked
souls, who always have the best hope, rose-colored expectations, always good omens. But their omens
are only the cries of birds, and their ``best hope'' like the sun in April. When expectation rides
summer into town, and the sunshine of the prospect flatters the wish, they are always delivered of a
new hope. But the newborn unfortunately gets only a short span of life: the prospects change, the
wind turns, and the blazing hope blows out. If, amid the changing prospects, we are to hold fast to
hope, then we must learn to hope.

The best teacher for memory and hope is fear; we must learn to fear. He who does not know what it is
to fear does not know either what it is to hope, and he who has no memory of danger experienced does
not yet rightly know what it is to remember. But precisely with respect to fear and danger it is
remarkable how we human beings apply a quite different measure to hope than to memory. The less the
danger seems to be, the greater is our hope; but conversely: the greater the danger has really been,
the greater is the corresponding memory. The threatening danger is a darkness of the future, a
possibility that awakens fear. Against the dark possibility Hope arms herself and points to the
bright possibilities, the possibility that the danger may pass by, the possibility that it can be
overcome. Hope looks askance at the danger, would most of all like to hide it, cover it with a flap
of her cloak, so that its terrifying appearance shall not frighten us. Hope enters into league with
the wish, and the wish and Hope agree that it would really be best of all if there were no danger at
all. The danger is the enemy. It is indeed possible that
% --- p. 6 ---
\setcounter{footnote}{0}%
\opage{6}the enemy flees when one boldly steps before his eyes, but the opposite is also possible. It
is possible that the enemy can be felled when one courageously gathers one's strength and strikes a
blow that hits; but alas, the opposite is indeed also possible. So fear alternates with hope, and
neither the wish nor hope rightly knows the truth about the danger.

Otherwise with memory; it knows the truth about the danger, for the danger has now been experienced.
Memory knows that a danger, seen at a distance, can seem smaller than it really is, and it teaches
us to fear. But it also knows that the dark possibilities which envelop the danger, flutter about the
danger like a traveling cloak as long as it still has the motion of something to come, can give it
a far more terrible appearance than reality itself, once it has come. The dark possibilities are
like thunderclouds; he who is in the midst of them is said not to be so exposed as those who move
beneath them.

There is especially one class of dangers which memory loves, indeed fosters with a fondness like one
who fosters flowers, dangers whose splendid pictures it always shows forth when it wants to coax the
wish and flatter hope; these are by no means those that passed by like storm clouds without a storm,
but precisely those which sober courage and the endurance of a resolute will happily overcame.
Well known is the word with which the warriors of antiquity encouraged one another when in earnest it
was a matter of staking their lives: \textit{olim hæc meminisse juvabit,} one day it shall be a
pleasure to remember this. What is bitter, heavy, indeed terrifying, when it is to be experienced,
can be consoling, beautiful and gladdening when it has been experienced. There is in poetry, in
history, indeed even in the
% --- p. 7 ---
\setcounter{footnote}{0}%
\opage{7}daily life examples enough which teach us what the memory of danger overcome can do, and to
what degree it can flatter hope.

It is told in a romance how King Alfonso of Castile --- after a brilliant victory over the Moors at
the river Salado --- sits on the throne and plucks the jewels from his crown in order to distribute
them as reward among the valiant knights. At last he has only one left, the ruby, the most precious
of them all; but where is the hero, the young hero, he who always kept at the King's side? In the
midst of a forest of Damascene blades his shield warded off the deadly blows; the King owes him his
life, where is ``the green knight''? Then the modest youth steps forth; but, he says with his hand
on his heart: the sparkling ruby is no reward for me. The jewel that alone rewards me, the King bears
in his breast. All are startled, angered, how strange! Then beneath the lifted visor a lovely
woman's face is bared. ``Donna Blanca, the Queen!'' A whisper went through the hall. That was a
moment. Such a moment, dedicated to memory, haloed by hope, yes, such a moment was surely worth
experiencing. How terrible to face the horrors of war when expectation threatens; but how uplifting
to have experienced even what is horrible, once the danger is overcome.

And now from daily life, a short story about a young mother. It is New Year's Eve; she sits alone,
her husband is on a journey. And yet, she is not alone: a little boy with clear eyes and yellow
curly hair is with her; but he certainly cannot think of keeping her company, he is so busy with
himself, he is drilling on the floor. As she sits in
% --- p. 8 ---
\setcounter{footnote}{0}%
\opage{8}thought and a quiet prayer passes through her mind, she feels her head heavy; she must lean
back a little and close her eyes. Then she sees a vision, a sickbed. In the bed lies a child, sallow
and wasted; a woman worn with weeping sits by it and stares at the little one, who opens his
half-glazed eyes and closes them faintly. Then the vision vanishes, she starts up in fright and
dread: was it an omen? She dares not speak of it at all. Her husband came back, and many a
confidential communication took place about what had been seen and had happened in the meantime;
but the riddling, the uncanny thing she could not touch upon. She struggled with fear and clung to
hope. So the year passed, and what then did the omen mean? The next New Year's Eve the little boy
had grown not a little bigger, and then also now and again a little noisier. It was a great joy to
the parents that they could see, that very evening, how he had grown in the last year and become
healthier and stronger than they had dared to hope. For quite truly a great change had taken place
in the little one. He had been sick, long sick, sick unto death. There had been, for the mother
especially, heavy days and dark nights. There had been watching at his bed, weeping at his bed,
praying at his bed, when all at once, just as the last spark of hope was about to be extinguished,
life wondrously returned. So the happy pair sat, husband and wife, together and explained to each
other how good it is after all that we do not know the future beforehand, how beneficent it is that
dispensations come upon us, that in man's mind there are powers hidden which first come to light
when the hour of decision is there, that we yet love most of all that for which we have most
suffered and striven, that which
% --- p. 9 ---
\setcounter{footnote}{0}%
\opage{9}we have won through fear and danger. So they sat hand in hand; it was mutual instruction:
they learned to remember and they learned to hope.

Yet we are not yet done with danger; it has more to tell us, when it is to speak out fully. Memory
has presented it in a softened light with the help of the joyful outcomes. The joyful outcome robs
danger of its sting, for the dangerous thing about danger is precisely the possibility of a
sorrowful outcome. Let memory now reckon it up, count the joyful outcomes together in one column, the
sorrowful in another, and see on which side it gets the larger sum. Even if both sums were equal,
indeed even if the calculus of probability found, what it does not find, that, on the assumption of
roughly equal possibilities, there must go ten joyful outcomes to each sorrowful one: it would still
not change the matter, not take the sting away, not remove the terrible possibility which is the
innermost core of danger. The possibility of the sorrowful outcome is the sword of Damocles that
hangs threateningly over the crown of every bold deed which is dedicated only to the temporal; all
enterprising endeavors, all hopes of the moment, all undertakings which fill the busy mind of the
race and saturate the present with content, have this threatening sword over them. No plan so
clever, it can after all founder; no will so stiff, it can after all be broken; no tower so lofty,
it can after all fall; no building so proud, it can after all collapse. Where danger shows its
Medusa head, there hope swoons, and memory writhes in pain.

So there must then be a power of existence which rules over the sorrowful as well as the joyful
outcomes. But what kind of power is it? The
% --- p. 10 ---
\setcounter{footnote}{0}%
\opage{10}whim of fortune, the course of fate, a guiding providence? What is it, who teaches us it,
where do we get certainty about it, a certainty that does not deceive? The origin of life, life
itself must give us certainty. But, it is said, life itself is indeed so very various, the one gets
one impression, the other another. Life speaks with a thousand tongues, each human being hears with
his own ear, translates into his own language and gets his own explanation. The one seeks the origin
of life in nature, the other in spirit; the one assigns the individual his goal within the boundary
of temporality: after death there is nothing to hope and nothing to fear; the other places the goal
in eternity: after death there is everything to hope and just for that reason everything to fear as
well; the one goes to meet danger with trust in the absolute immutability of the laws of the world;
the other believes in a fatherly providence, a love which never sleeps. Who, now, is right? It is an
old quarrel, and the question is how it is to be settled.

It might well seem that, if the solution of the riddle were to be sought in life itself, the many
views of life could not possibly give such contradictory results; it might seem that the scientific,
calm, sober, thorough, comprehensive view in particular must at last become able to master the task
and settle the dispute. Herein lies an error. Life must indeed have a foundation unchangeable in
itself, but how deep it lies and of what nature it is, is discovered not by brooding, but by living:
the solution of the riddle is not in the outward, but in the inward. The inward is not exhausted in
the universal, which all have in common by necessity; there are springs at the bottom, and each
individual must, on his own account, see to digging and finding. Human life
% --- p. 11 ---
\setcounter{footnote}{0}%
\opage{11}certainly takes on an ambiguous appearance when one regards it outwardly, wants to survey it
as one surveys a landscape, or to follow its course like a river whose source and mouth it is a
matter of getting accurately marked on the map once and for all. If it takes long to discover the
sources of the Nile, it will take longer still before inquiring science gets so far that it really
discovers the sources of life. Even if the discovery is given in a scientific silhouette, not even
the half is thereby achieved. If we call life a river, then this river differs from the other
rivers of the earth in that its outlet is veiled in an impenetrable darkness. It loses itself, it is
said, under the earth. Yes, under the earth! What then is it in its further course, a Lethe, muddy
and sluggish, the river of oblivion with moldering bones and the rubble of crumbled bauta stones? or
perhaps a bath of rejuvenation, a mirror-clear water with heaven's depth, a river of paradise, that
stream of hope from which memory drinks?

The final goal of life cannot be discovered by any outward contemplation whatsoever, cannot be fixed
by a thinking consideration of the universal, be it of the stars of heaven or of the flowers of the
field, be it of the universal life under the forms of necessity or under those of freedom. Every
consideration, even the most ingenious, of the appearances of the self-being in nature or history
is burdened with ambiguity and gives us to every question a double answer, both yes and no. Hereby
neither the use of reason nor the aids of experience are rejected. The laws of nature, the destinies
of mankind, the course of history, all contains hints and guidance. We honor the traditions in the
Church's teaching; the
% --- p. 12 ---
\setcounter{footnote}{0}%
\opage{12}holy words of revelation are to the believer a rock foundation on which he builds his hope.
But that does not change the fundamental law of existence, not the conditions of life; for this is
the condition of life: that everyone who wants to get to the bottom of the great question must labor
to get to the bottom of himself. The ancients believed that somewhere in the earth there was a cave
through which the way went down to the underworld, down to the realm of the dead. In such a cave we
do not believe, but we are assured that there is in the soul of man a depth of which many have
scarcely more than an inkling, that it is the dim depth in which eternity dwells, and where the way
goes down to the land of riddles.

That this conception of the inward agrees exactly with what shows itself in a consideration of the
outward, history can teach us. If we go back three hundred years, what a difference between past and
present, precisely with respect to the question of life and hope!

The question which drove the Reformation forward was, despite all temporal conditions and
admixtures, certainly a question of eternity, and to that extent a mighty, decisive question of
life; but it was in a form which the critically reflecting present will not acknowledge. In the age
of the Reformation and of the Thirty Years' War it stood fast as a settled matter that positive,
revealed religion must be the foundation of life and hope. About that there was no dispute, no
doubt. But in which church community is the unadulterated religion to be found? About that there was
doubt and dispute and bloody strife. For the present time the first decisive question is not: in
which church community has positive religion its purest expression? The question is much
% --- p. 13 ---
\setcounter{footnote}{0}%
\opage{13}more this: does human society really need religion? Can human life not complete its course and
reach its goal without the assistance of religion, at least without that of positive religion? In
those days it was a fixed dogmatic presupposition that the soul must continue its individual life
after death; the question was: how will it continue it, in blessedness or in damnation? Is there a
purgatory, or is the doctrine of purgatory a fable? In our day the question of purgatory and eternal
damnation does not stand in the foreground; the question is whether there is a life after death at
all, whether the hope of immortality is not itself an illusion.


Seen superficially, this change looks like a sorry step backward, not like an essential advance in
the history of culture. It must indeed seem as though the dogmatizing fathers in the sixteenth and
seventeenth centuries took the great questions of life in greater depth and with far greater
earnestness than their politicizing sons in the eighteenth and nineteenth. From the standpoint of
individuals it may perhaps be so, but not from that of history. The general consciousness has, as
regards the religious, made an unmistakable advance in our time. The presuppositions on which those
dogmatizing fathers built their faith and hope were all supplied from outside, all passed into
crystallized form and stiffened into outward authority. This holds of the Evangelicals just as well
as of the Catholics. Every doubt was branded as heresy, every contradiction struck to the ground by an
authoritative decree, everyone who thought otherwise was instantly forfeit to judgment. The
consequence was that the deeper life of the soul, despite stirrings, longings, sighs and tears,
% --- p. 14 ---
\setcounter{footnote}{0}%
\opage{14}never attained full development or clear understanding. Faith was strong when passion flared
up, but just because the believing gaze one-sidedly fastened on one outward, fixed determinacy of
form to the exclusion of every other, enthusiasm turned into fanaticism. The most powerful
nourishment of fanaticism is hatred and damnation. In our day fanaticism has given way to sobriety,
the spokesman of damnation has grown weak of voice. The consequence of this is that the inward can
come to free development, that it can become recognizable what is stirring in souls. Different
voices make themselves heard, opposed views clash with one another; but no authoritative decree
breaks through immediately, no outward authority gets the right to decide a dispute which can only
be fought out with spiritual weapons. Faith is impossible without inner freedom; but freedom for
faith is at the same time freedom for unbelief. Only on the condition of freedom may it be hoped, as
the struggle proceeds and the crises become more serious, that the old quarrel, the long lawsuit,
may at last be brought before its proper court and the case be referred for judgment.

That human life, according to its innermost essence, must be of one design, have one common goal,
and that the truth must to that extent be the same for all, cannot be disputed. The mutually
conflicting views of life arise from a superficial, one-sided conception, or even a denial, of the
fundamental truth, and consequently from a misapprehension of the \emph{original} essence of life,
its way and goal. This misapprehension cannot be removed by norms of truth set up outwardly. Let it
be a thousand times true that the truth has revealed itself once and for all; what does that help,
when its manifestation
% --- p. 15 ---
\setcounter{footnote}{0}%
\opage{15}is misinterpreted. If it were an immediate truth of nature that was at stake, it would have
to be confirmed by investigations of natural science; if it were a plain historical truth of fact,
it would have to be established by historical criticism; if it were a general truth of essence and
necessity, it would have to be tested in objective philosophical thinking; but a personal truth, a
truth of will, a truth of life grounded in freedom, how is it to be tested? What freedom is can be
recognized only by one who is set free and appropriates freedom; what personality is, only by one
who has dug so deep into the ground of the mind that he has reached the origin of life. Everything
depends on experience. The many mutually conflicting views of life thus arise from the many
different, lower and higher, stages at which experience stands. If human life in its wholeness has
one design, one fundamental law, one goal, then it is impossible that the demands of life can be
satisfied equally completely at any stage of experience whatsoever.

That someone lives either at an immediate or at an arbitrary stage, in immediate enjoyments or
arbitrary interests, and finds himself content, is no reliable mark of truth. It may, after all, be
good fortune, accidental individuality, a condition determined by circumstances. When the
circumstances change, and fortune fails, what then becomes of the accidental individuality? Only he
who has a certainty of experience that neither on earth nor in heaven, neither in time nor in
eternity, is there any power of existence that can bring him into discord with himself and lead him
away from his goal, has recognized his origin and knows what life is. The only test of a view of
life is this,
% --- p. 16 ---
\setcounter{footnote}{0}%
\opage{16}that life is allowed to unfold itself and gain shape. If there has been sin against the
fundamental truth, existence will surely make it manifest, and the error will come to light.

It is decidedly certain that a living being can never become anything greater or better than that
for which it is originally laid out; for man it is a matter of knowing the design of his being. But
man's design is double: a fundamentally essential one, which corresponds to the race, and an
individual one, which shows itself in a particular endowment. It is important in both respects not
to rate oneself too high and not too low either. As regards the individual endowment, the individual
may perhaps from the beginning rather be tempted to expect too much for himself than too little, for
imagination is in motion: hope with her magic mirror goes ahead, and on the way there are nothing but
roses; but when (practical) experience comes with experience, the thorns pierce through, and it may
then happen that memory gets scars and scratches. As regards the general endowment: being a human
being, the race, particularly in our day, when natural history flourishes, is rather tempted to
rate its design for life too low than too high. To be endowed as plain and simple man has nothing
flattering for the individual, but to be a talent, a genius, outstandingly endowed, can flatter the
individual. When someone says: we are an animal species, no doubt modified apes, he does not thereby
mean to offend anyone, least of all himself: he is, after all, merely bringing his knowledge of
natural history to bear. But if someone is present who, with old-fashioned breeding, would answer the
speaker: then you yourself, on the score of your descent, are at bottom a brute, that would be an
insult. To be an animal, when we are all animals,
% --- p. 17 ---
\setcounter{footnote}{0}%
\opage{17}is not insulting, it concerns only our common endowment; but to be called a brute is
insulting, for it is to be charged with a lack of individual endowment.

We ought not to give ourselves over to fanciful expectations and set our existence too high; but
neither ought we to give up hope, the hope of eternity which ennobles our race, we ought not to set
our life too low. There is in experience a measure for hope, which memory keeps. It is this measure
that we shall bring forth in what follows, as we speak of life and hope.

% --- p. 18 ---
\setcounter{footnote}{0}%
\lecture{Second Lecture.}\opage{18}There is a legend of a wonder-working spring, a
well of health, healing for all who are weighed down by
the burden of the years; it is called the fountain of youth. It is natural
that in all the regions of the world, in our old
Thule too, diligent search has been made for such a spring.
But all investigations have hitherto been fruitless.
The poet with his wishing-rod has untiringly gone at the
head of the seekers. Often the rod has bent in
his hand, and he has given notice now here, now there;
but alas, each time one rushed up and stuck
a finger into the soil of reality, it proved to
be --- a mistake. Archaeology too has in
a way taken up the matter by undertaking valuable
excavations; but the archaeological congress has, as far as is known,
probably not yet come to a clear conclusion
about the result. Modern criticism, on the other hand, has long since
come to the clear conclusion that the fountain of youth must be reckoned
among the lost causes, that it is with it
as with Paradise, the Philosophers' Stone, the Garden of Irem, King
Solomon's Ring and other glories of the past, which are
lost. A pity, should it be true that the spring is not
% --- p. 19 ---
\setcounter{footnote}{0}%
\opage{19}to be found. We dare not give up hope, and when we
do not give up hope, we find at once what we seek; for
the fountain of youth is, after all, precisely hope.

That hope is a fountain of youth may be seen already
from the fact that there is always something youthful in
hoping, and that it is above all the young who hope; but
the matter first receives its right understanding when we
reflect that there are two kinds of old age, a
bodily and one of the soul, and, corresponding to this, two
kinds of youth as well. Against the old age that comes with
astronomical years and wants to be reckoned by the almanac, no
herb grows; for it there is no fountain of youth. But
the soul is not thus dependent on time; it knows
a youth that is already old, and an old age
that is still young: there are old men of eighty
who are nonetheless younger than men of thirty. When, in
the sorrowful sense, did a human being grow old? On the day
it stopped hoping.

As we speak of two kinds of old age and two kinds of
youth, we must above all not forget that there are also
two kinds of hope. There are temporal hopes, and the
temporal hopes are many; and there is a hope of eternity,
which is always one and the same. The two kinds of hope
can be separated. One can hope immediately for the
temporal in forgetfulness of the eternal, and there are
those who think that the heart's dearest hope first truly thrives
and blossoms in the sun of the passing day, when the cool breath
of eternity is gone. One can hope for the eternal with
abandonment and renunciation of the temporal, and there are
those who set the enmity between temporal and eternal
so high that the earth must become a vale of tears before we
get so far as to descry heaven. In
% --- p. 20 ---
\setcounter{footnote}{0}%
\opage{20}reality, however, putting off the temporal does not go so simply
as zealous preachers of repentance seem to suppose; it
is an intricate divorce case. It is eternally certain that the
two kinds of hope must not be confused; but it is also
certain that in this present they belong together as
spirit and nature, as soul and body. What is a
temporal hope when the very last gleam of eternity,
the last sunlit glimpse of ideality, has vanished? It is
no longer hope, but either a selfishly covetous
fancy or a cold expectation founded on probability.
And how does hope look when
it turns away with contempt from every content
that offers itself through the temporal? It is an
experience of life that one not only goes wrong when either
the so-called natural is to be made one and all, or
one lives on a false compromise between temporal and eternal,
but that one also goes wrong, and that the error is no
less dangerous, just because one aims higher, when the ground of
natural life is to be cleared away so that the spiritual can
find room. It is true, to be sure, that the multitude of men,
especially in our century, is far less exposed
to the latter temptation than to the former; but nonetheless
we ought to be on our guard against the most dangerous
of all errors. It is no rare sight in the history of religion
that worldly lust and pious mortification
go hand in hand, that the selfish nature can just as
well be brought to rage in renunciations as to
wallow in enjoyments.

The true task first presents itself when one
has seen that the natural belongs to the development of
spirit, that the eternal borrows material and content from the
temporal. From this standpoint it becomes clear to us
% --- p. 21 ---
\setcounter{footnote}{0}%
\opage{21}what the youthful means, how youth in the
essential sense is related to hope.

In the youthful mind two elements, the
natural and the spiritual, or, what says the
same, the real and the ideal, the temporal and
the eternal, are originally grown together; they are of one shoot,
branch and blossom of one bud. The young, inexperienced,
but still uncorrupted soul hopes, and for what?
Naturally for all that is noble and great and good
and beautiful in this great, beautiful world. It is not
the empty concept of greatness, not a general idea of
a general good or a reflected-up notion
of the in-itself indeterminable beautiful, that fills the hope of youth.
No, the great has shape: heroes with swords;
the good has flesh and blood; and beauty? Why, she sits
at the window pane, if one only dared to greet her.

It is a mistake when one would overfill a youthful
imagination with tales of knightly tournaments and
prizes of beauty, of heroic strength that never fails, and
of the many wreaths of virtue which the fingers of Fortune
are busy weaving for the victors. It is an
instruction that puffs up; one gives the young overstrained
notions, expectations without experiences,
bombast and sponge, and hatches nothing but fantasts. But it
is more than a mistake, it is a crime against
the young souls, when instruction is systematically arranged
to make them old before their time; when the
mirror of reality that is constantly held before their eyes only
shows them images of an outlived, perished world
with false ideals, a theater of civilization where
impudence and intrigues and lust for power and hidden despair
play the leading part; when memory day by day
% --- p. 22 ---
\setcounter{footnote}{0}%
\opage{22}is loaded with fragments from the glory of the devastations,
while the understanding busies itself with forms and forms of
forms without end. The owl of wisdom builds in the ruins of
greatness! But when no word of life, no ray of
spirit lights up this underworld of sighing Titans
and fallen Giants --- only connection and web, and
web and connection is the whole explanation: what
wonder if the young man unconsciously, half-consciously gives up his
first childlike hope; what wonder if he little by
little, even before the time of experience, makes himself familiar with
the course of fate and the play of forces and the superiority of the calculating
understanding; what wonder if a coming
generation experiences that old men of twenty shrug their
shoulders with a smile if there is still a fool so
youthful as to want to swim against the stream of probabilities
and stake his life on truth and right. There is
one more crime of the spirit, which steals youth from
the young.

If the hope of youth receives the necessary conditions
for a healthy and natural growth, the relation between the
temporal and the eternal will work itself out in the course
of life, and the direction will differ according as the character
develops. If a character is really to be formed,
and the demand of personal freedom to be satisfied,
then a turning point must necessarily come,
when the will chooses one of two things: either to posit
the temporal as the only real and deliberately
exclude the eternal, or: to posit the eternal as goal
and work up the temporal as an element of development, as
means. How this happens in the lives of individuals cannot
be described: here there are as many possibilities as
individualities; but the fundamental law by which the development
% --- p. 23 ---
\setcounter{footnote}{0}%
\opage{23}proceeds can be shown, and that is what I shall now
attempt.

Hope is determined by what is hoped for; but as
long as the temporal is still grown together with the
eternal, hope finds itself, so to speak, in a pupa-like
state. The transformation takes place imperceptibly. So,
e.g., with the hope of love. Those who hope hardly know
of it. What has been settled between them in happy understanding
is something temporal and yet --- for ever! Then
comes the moment: hope flutters like a butterfly
from flower to flower, but the happily hoping feel
assured that the light wings can bear it
up to the stars. They have, in other words, a presentiment
that an eternal life can begin in time. The latter
is right; but the understanding is so youthful. The
youthful hope must be tested; in what, now, does
the test consist?

The young hope is prepared for obstacles, ready
to fight against resistance. If resistance and
obstacle were not possible, it would not be hope,
but mathematical certainty. Resistance is not hope's
most dangerous enemy, for hope has courage; the obstacle, even
if it stretches far, is not hope's most dangerous
enemy, for hope has patience too; there is a perseverance
in hope that whispers: give it time, my day
will come, fortune shall crown the faithful. But there
is a possibility that time holds in store, an enemy
that it treacherously harbors: it is disappointment in the form of deception.
Deception is hope's most dangerous enemy; it does not postpone
fulfillment, ah no, it brings fulfillment
to nothing, it would assassinate hope. But whence
comes deception? Not from the eternal, no, from the
% --- p. 24 ---
\setcounter{footnote}{0}%
\opage{24}temporal! Here the two elements, eternal and
temporal, begin to separate, and this separation is a crisis.

It would be very youthful if we supposed
that it was only the hope of love that was thus exposed
to deception. The possibility of deception lies in wait everywhere;
if there are manifold temporal hopes, there are also
many kinds of deception and manifold disappointments. A
disappointment is painful in the same degree as the hope was
trusting and rich and beautiful. The most dangerous thing about the crisis
is that the disappointment of a single hope is so easily able
to spread, to attack hope itself and make the
disappointed one hopeless. A solar eclipse, where it is
total, spreads a deathly pallor over the whole of nature;
the beasts of the field are seized with dread, the cattle bellow, the small birds
flutter about bewildered or make for their nests.
So also, when the shadow of deception covers hope
completely, the view of the whole of life will undergo a
transformation. It is not merely this or that single
prospect that is darkened, but the foundation of trust is shaken,
existence takes on the gleam of a sepulchral lamp. Why
do we live? Only to die. Why do we enjoy? The better
to suffer. Why do we rejoice? To grieve deeply. Why
do we hope? To be bitterly disappointed.

As long as the mind is in turmoil, the consolations
of those around will naturally find no entry. That it
will surely pass, that time heals all wounds, that by way of
compensation one can of course get oneself a new hope, such everyday phrases
are rejected with contempt: ``a new hope, yes, with
a new disappointment, a possibility which in time, when
Fortune one day is in her evil humor, is allowed to
set deception loose!'' We note that the crisis here sets up
a divorce, not merely between past and future, but in
% --- p. 25 ---
\setcounter{footnote}{0}%
\opage{25}the strictest sense between temporal and eternal. Whether the
separated elements are ever in this life to be reunited,
and how, depends on more than outward
circumstances; it depends on the resilience that
is in the mind, and above all on what the will chooses
and the mature resolution holds fast.

There are probably few human beings who have trodden
the path of experiences without having also in one way or another,
with dimmer or clearer knowledge, in weaker or
stronger degree, experienced what it means
to go through a crisis of disappointment. We cannot simply
fetch our examples from the outward, least of all
from public life, for matters of this nature do not go
on exhibition. History preserves them in secret
archives. But if we begin from within and consult
psychology, it will supply specimens in
abundance. Moving groups of individualities, men
and women, youths and girls, poor, rich, simple,
clever, present themselves here on this first, dangerous
stage of transition. Let us observe them more closely: figures
on a dark ground, physiognomies of an apparently
confident exterior, but the glancing light falls on an enigmatic
interior; it is as if each individual had a double
face. They are all sufferers, although one cannot quite
see it on them. A little pale, perhaps; but paleness
can, after all, have many causes. Some thoughtful,
brow clouded, eyes swimming; others who raise
their heads proudly, in whose smoldering glance and clenched
lips one as it were reads a challenge; in short: from the
exterior nothing can be inferred. What gnaws at them is a
secret; they are all wounded by the sting of disappointment,
all bitten by the serpent of deception. Young, and yet already
% --- p. 26 ---
\setcounter{footnote}{0}%
\opage{26}old! For as one knows that the hair can turn gray in
a single night, so a disappointed hope can in an hour
of grief make a human being older than time with its calm
pace through many years. Yet the illness is not
mortal. It will pass; the wound fever subsides,
healing powers work in hidden workshops, in
a short while life again asserts its claims: we can
soon regard our sufferers as convalescents.

The point on which the psychological, we may
add the religious-ethical, interest rests is now this:
how will the healing take place? With those who,
without grasping the ambiguity from which the crisis was to
tear them, get themselves a new hope as one gets oneself
a new dress, and begin afresh with something temporal
and yet for ever, we are easily done: they were not
really wounded, they had only fainted. We can then
divide our invalids into two classes. To the first class
we reckon those who with a certain resolute coldness
of resignation take hold of the real with abandonment
of the ideal, choose the temporal with their backs turned
to the eternal. They are in the most proper sense without
hope and yet not hopeless; for he who voluntarily dismisses
hope is not hopeless. The temporal has
many claims and can therefore offer them many expectations.
Only they expect no rejuvenation; their
motto is: no one deceives us, for we are old.
Under the second class we bring those in whom the
eternal, the ideal, was only pushed back in order to react
from the bottom of the heart with still greater force. They drink
anew from the well of life: the fountain of youth trickles, hope
grows and is transformed. It goes without saying that each
of the two classes must in turn fall into various
% --- p. 27 ---
\setcounter{footnote}{0}%
\opage{27}subdivisions, according as the mind is weaker or stronger,
according as it is feeling or understanding, imagination
or will, that predominates.

First a couple of words about the first class. The weaker
mind, which has lost its hope for all that is higher,
sinks down either into the petty occupations of everyday or into
the stupefying tumult of the passing day. It learns from experience to
adjust its expectations to the circumstances;
but even if the outward conditions are brilliant, even
if sensual passions flare up and find satisfaction,
there is nevertheless an emptiness, a quiet bitterness in
such a soul, a gnawing despondency that the world is false,
existence hollow, that life does not keep what it promises.
Individuals of this category may, for that matter, be
good-natured, rejoice with them that rejoice and regard the happiness of youth
with a peculiar melancholy of sympathy; but
they may also be envious, malicious, and
measure those who hope with a look of malicious joy that says:
you too shall be disappointed!

It is otherwise with the strong mind when it has taken
leave of hope, exchanged its ideal aims for
great interests of reality and devoted its strength of will to
serving under the present. That such a character
should become a bad or dangerous human being is by no means
given. It can become a man honorable in his circle,
an enterprising man, quite a man of the times, a businessman
from the top drawer. There is nothing to prevent such a
man from having a benevolent mind, a good heart, from
being able to rejoice in working for the benefit not only
of himself but also of society. But he has
a view of life that he does not preach in the street; it
is imprinted on his mind so deeply that he does not need
% --- p. 28 ---
\setcounter{footnote}{0}%
\opage{28}to have it confirmed by others; its short wisdom
is: when the coffin lid is shut, the story is over. With
this worldly wisdom he does well enough under the circumstances.
Certainly it would be insufficient if
death really should not be able to cut the
knot in Alexander's manner; but that is death's
affair. A man in such a business has no time to
brood over such things; he takes it practically. He meets
adversity, he experiences grief and sorrow, but knows how
to adjust himself to everything. It is as if with every
year he got a new confirmation of the truth of his view.
Quite naturally; for the element in which the soul lives,
and from which it draws its nourishment, acts upon its
need and nature. One feeds young animals, e.g. rabbits,
with madder in order to see the red layers in the bones. A
similar experiment can be made with the soul. Man's
soul gets temporal juices in its marrow when it is daily
sated with the temporal, only with the temporal. The more a
soul gives itself up to impressions of the sensuous, the more
easily it convinces itself that nothing has reality
except what is sensuous; the deeper an individuality
with all its interests sinks into the finite,
breathes and lives in sheer finitude, the more evident
it becomes to it that the eternal is a semblance, the beyond
a dreamt ideal: one can involuntarily perceive in
oneself that the soul dies with the body. But one does not on that account
despise virtue; on the contrary! temporal life has
need of virtue, and therefore one must practice virtue. So
this virtuous man lives in restless activity as
father of a family, citizen, politician, at the head of manifold
public undertakings. Not only joint-stock companies
and charitable foundations, the Church too owes him much for important
% --- p. 29 ---
\setcounter{footnote}{0}%
\opage{29}economic improvements. Thus
he lives, and so he dies. And he gets a funeral
--- three long streets full of carriages --- and a
funeral sermon, yes, a funeral sermon! for now this
man of honor goes in to eternal blessedness, into the regions
of transfiguration, to fetch his reward, a reward so great, so
well deserved, so blessed and rich, that there is not a dry eye
in the church. --- How this man has understood
himself, what responsibility he incurred on the day when,
wounded by disappointment, he offered his hope of eternity for sale against
a handsome sum in temporal currency, is a matter by
itself. But in one thing he has understood himself; for we know
him to the bottom --- we have not met him in the outward
world, but found him in psychology; we have,
after all, created him ourselves, and so must surely know him
--- he has in his way rightly understood that he and
eternal blessedness have not the least to do
with each other. Could he answer from the coffin in there,
a voice would be heard: I got the reward I expected;
hush about the eternal; let me now have peace in my grave.

If we turn to the second class, to those who
drink of the fountain of youth, the individual difference
is again to be recognized according as the new
hope remains in mood and imagination or makes a covenant
with the will and steps out into character and action.
It is the rebirth of hope that we are to consider. How
can hope be reborn? It dies in disappointment;
if it is to come to life again, it must be born anew.
What we in this connection call the reborn
hope is not new in the sense that it should have
another origin and be quite other than the one
with which we began. The hope of youth was precisely
% --- p. 30 ---
\setcounter{footnote}{0}%
\opage{30}so beautiful for the very reason that it united in itself the eternal
with the temporal. It was the crisis that set discord
between the two elements. The reborn hope must then
take up the temporal once more, but, it goes without saying, on other
and higher terms.

We must constantly remember that life itself is the spring
of hope. If life has from its origin a double claim, a
temporal and an eternal, then hope too must have a
double content. If hope did not have an element of eternity
in its content, it would, as we have already
seen, have to go under in temporal expectation, give up its
ideal possibilities and lose itself in a calculation of probabilities.
If hope dared not borrow a content from the
temporal, it would have to exclude unconditionally what is on this side;
like an outlaw, an exile on earth, it would have to
stare unceasingly toward the beyond and seek its home
in the blue. But the blue and the eternal are not one
and the same. If hope is to show the way through life,
it must cast its glance on the real earth: he who
walks on the earth with eyes always turned heavenward easily exposes
himself to stumbling. But when hope again
grasps the temporal, the one who hopes, having come through
the crisis of disappointment, must be well aware that the
temporal is grasped for the sake of the eternal. What matters
unconditionally is an assurance that the
ideal itself is the ground of hope, that hope is fulfilled in
the eternal sense each time it gives us a confirmation
of the reality of the eternal. On this presupposition
hope can in this present life be both disappointed temporally
and yet fulfilled eternally. Indeed, the temporal disappointment
can, as we shall see at once, be a powerful means
of confirming an eternal fulfillment.

% --- p. 31 ---
\setcounter{footnote}{0}%
\opage{31}It is overexertion when a modern wisdom
undertakes to draw heaven down onto earth and to recast
existence, so that the transfigured moments of the collective life
should make the eternal itself commensurable
with the temporal: we are not yet in heaven because
the earth is located in the space of the heavens. But it is
no less, for all that, a dangerous misdirection, ruinous for the life of the spirit,
when the eternal is made into a beyond in the sense
that the temporal becomes sole ruler on
earth. The truth is that the eternal, with inner independence,
from out of its own kingdom, works through the
temporal, bears the temporal, strives to reshape,
transfigure and transform the temporal.

Let us, so as not to skip too many intermediate links,
begin with imagination and mood. Think
of deeper, fermenting natures such as the mystics; think
of individualities of genius in whom the ideas are at work.
They are by no means souls who all hope for
a personal life after death; and yet they have an involuntary
sense of the eternal. It dwells in them,
and precisely because it is so immediately present to them,
they can --- without there yet being any question of a beyond
be an example to us of what the rebirth of hope,
what we call its eternal fulfillment, means even in negative
form. Take good note of the poets,
especially the restless ones, those whose innermost being is pain:
Byron, Heine, Alfred de Musset and many among the
very newest. However much they praise the natural side of our life,
the blind but infallible instinct and the sensual,
elemental passions, however much they deify
the female body in particular: one notices nonetheless that
they are just as dissatisfied with fading nature and
% --- p. 32 ---
\setcounter{footnote}{0}%
\opage{32}its withered flowers as they are dissatisfied with the
philistine emptiness in the spirit, with the conceited,
petty mind of men and the staid prose of society.
What is it that, in the midst of brilliant company,
despite the laurels of recognition and the incense of admiration,
makes them lonely birds? What kind of
gleam of light is it that in one and the same instant delights and pains
the eye? The sparkling ideal, on whose lightning flash they set their hope!
Constantly the ideal demand and then a disappointment. Upon
each disappointment follows a new flash with a new hope and
then --- again a disappointment. Hope dies in disappointment
to be born anew: here is an element of eternity that
unceasingly reacts against the temporal. What torment, that
hope in instant after instant both dies and lives! Unhappy
throes, a process of rebirth with a hopeless hope, and
why? Because it stays with mood and imagination
in a selfish spirit, because the eternal never gets the strength for
real redemption.

If hope is to be reborn in character and will, other
figures must come forward on the stage. Here is a
man; we do not fetch him from the street, he comes with
greetings from psychology and ethics, ready to submit
to a test of reality. He is a human being
whom the pain of disappointment has taught that one should not selfishly
live for oneself, but win one's life by staking
it and working for ideal aims. Real life
has aims enough; it also has ideals. The ideal of reality
that shines nearest for him is freedom, not
the freedom that lets the multitude loose, but the freedom that
ennobles the multitude and lifts it up into a people. Here is
his goal. So he sounds the alarm in the sleepy camp;
his speech rolls like a mountain torrent, the masses awake,
% --- p. 33 ---
\setcounter{footnote}{0}%
\opage{33}crowds flock around the enthusiast and devour the
awakening words. It is a vision of the future with a great
hope that here unfurls the banner; will the hope be fulfilled?
Prudence knows: it will end in a disappointment.
Yes, alas, in a great disappointment. That is nothing
new, then; but what is peculiar to this man
is the way in which he bears himself in
the disappointment. He must surely, people say, grow wise by experience.
Yes, certainly he grows wise, but in his own way: only now
does he begin to act like a fool. Many
who in the time of awakening carried the banner high, went aloft
with the new flag and were bold skippers steering
straight for the maelstrom, cowered in fright when
the whirlpool came and the deep began to boil; wise,
taught by experience, they quietly went about, carefully
furled the banner and left ``the sinking wreck''
with the ideal in their pocket. But the man of hope held out
at his post; that was a youthful folly. And
one thing more, which was more youthful: he stood by
his word, he was still glad in his first expectation.
No one, not even the most cautious among the cautious,
had a scout's eye like his when it came to discovering
the danger. But the danger for him was this alone, that unbridled
excesses should choke freedom, and the multitude itself
lay the people waste. No one, not even the bravest among
the brave, dared like him so openly, freely and unsparingly
to point out the guilty and tear the mask off the seducers;
that, after all, was a folly. No one like him dared to lay bare
the gaping abyss into which greedy craving for power, in
league with selfishly divisive interests and a sallow,
sour-eyed, all-leveling envy, was about
to hurl the bewildered mob. Folly! He wanted,
% --- p. 34 ---
\setcounter{footnote}{0}%
\opage{34}not by mockery and scorn, but by honest, unvarnished speech,
by the clear, convincing words of truth, to open the
eyes of the blinded and lead the movement into the right track:
youthful folly. At last, when Nemesis came,
the goddess, the inexorable one, who punishes all excesses
and avenges reason that has been scorned, he stood quite alone.
--- Is he now without hope? By no means. His hope
stands fast, undarkened, and shines like a beacon from the mountain.
No one can understand that. It must surely be, people say,
because he despises the counsel of experience that he
can go on being so youthful. No, no one can
understand it. Clever people, who otherwise understand
everything, cannot understand him, for in his
experiences lies a secret. His hope has
perhaps disappointed others, but it has not disappointed him
himself. How could the hope of freedom disappoint him who
at every outward defeat had experienced that he
won an inward victory. How could the faithful hope fail him,
the hope that made him eternally free and pressed
the seal of liberation on every self-sacrificing action. He
did not stare up at the blue; he dug deep in
the earth of temporality, and he found a treasure. And this
treasure? It was more than a horn of gold, it
was more than the Philosophers' Stone, it was more than King
Solomon's Ring; it was a pledge of immortality, a spring
welling up at the bottom of the soul, a way of hope to the well of life:
he drank of the fountain of youth.

% --- p. 35 ---
\setcounter{footnote}{0}%
\lecture{Third Lecture.}\opage{35}Hope is of personal origin. Among all creatures
on earth man is the only one who can in the proper sense be said to hope; for man
is a personal being. The animal can desire, long, expect; but because it expects a good,
this expectation is still not yet hope. What ennobles an expectation into hope, even where
the temporal stands in the foreground, is a freely projected image of the future, a
possibility clothed in bright colors, a spiritual clarity, a glow of ideality, which the
animal's eye never sees.

Religion speaks of the whole creation's longing, its pain and sighing, its hope of deliverance
from the bondage of corruption. Romantic poetry repeats the same; it listens to a quiet
lament, a soft, questioning whisper, a melancholy sigh from the soul of the universal life in
nature. ``The red of evening, when it writes its good-night in the banks of cloud, asks whether
it shall not soon be the herald of a blessed day''. How is this to be understood? Who is it that
sighs and laments, who is it that longs and hopes? Not mute nature, but speaking man. It is
man who lends nature his voice. There are indeed voices in nature,
% --- p. 36 ---
\setcounter{footnote}{0}%
\opage{36}many voices, and yet nature itself is like a deaf-mute. The elements become loud when
they roar and clamor; but roaring and clamor are not yet a speaking voice, any more than the
deaf-mute's mumbling. The animal has a voice and its soul a sigh, and voice and sigh can tell
us much when they are interpreted; but the animal is not its own interpreter: it is a dumb
creature. It is man who lends nature his word; it is man's soul that makes itself the
interpreter of its pain, its sigh and its lament. And with well-founded right. Only in man,
the thinking being, does the opposition between the actual and the ideal come forth from its
hiding place; only in man, the self-conscious being, can natural life sigh under the bondage
of corruption; only in man, the personal being, can the whole creation hope.

As hope is, so also is its opposite, fear. The animal's fear is obscure presentiment,
involuntary fright and dread. Man's fear goes further, he has a personal fear, a fear of
recognized possibilities; the self-conscious imagination discovers the alarming possibilities
everywhere, lures them to itself, dresses them up and gives them form and shape. There is a
shape, dark and threatening, that stands on the road, now farther off, now nearer, always the
darker, the more sunlit the road is; it is death. The animal dies, but it has never seen death.
Man's soul sees death many times before it actually comes. Every time someone has set himself
a goal in something temporal and really means to look forward to its attainment, the shadow of
death stands there, and a struggle arises between the shadow and hope. You may, after all, be
dead before tomorrow! how disheartening for one who expects a joy, a truly great, unspeakable
joy, from the
% --- p. 37 ---
\setcounter{footnote}{0}%
\opage{37}day of tomorrow. Hence despair's grasping after the moment of enjoyment. Religion has
lent it a word: let us eat and drink; for tomorrow we shall die.

Even those who want more than to eat and drink, who want to work for higher purposes but yet
seek life's completion within the temporal, meet the shadow of death on the road; but courage
grows, the shadow must yield, the stout-hearted conjure it down with a powerful formula, and
the formula says: \emph{The individual dies, but the race lives}. There is something rousing,
enlivening in this thought, it has something consoling about it; let us weigh the impression
and understand what it says, in an image. --- There stood an old tree, bathed in the spring by
waves of warm air; the sun looked down glowing from heaven and preached to field and forest the
gospel of the resurrection with such pith and power that everything swarmed up out of the
earth. The buds on the old tree hurried to see who could be first to burst the mummy-like husks
and display the fresh, sappy shoots: there were already leaves and blossoms on the tree. But
some sunbeams, which played among the branches and peeped curiously in through every crack,
took notice of a single bud that sat so reluctantly, huddled in its husk and would not come
forth with anything. What is the meaning of this? they asked, why do you hold yourself back
like that? Oh, it answered, I am only an insignificant individual, there are buds enough
besides me. Besides, I cannot be bothered to burst out, it cannot lead to anything anyway. ---
What is that you say? Can it not lead to anything? --- No, it answered, I have heard a whisper
that it is only vanity and short-lived glory, all this bursting out; in a little while the
blossoms will fall off, and the leaves
% --- p. 38 ---
\setcounter{footnote}{0}%
\opage{38}be whirled about on the ground. --- That is wind, said the sunbeams, it is a wind that
has frightened you; do not listen to the wind, it always goes about with idle talk. Just you
come out, small as you are, and you shall see the great radiant world, and you shall see of
what mighty, glorious tree you yourself are a member. You shall be proud to make your modest
contribution to the tree's growth and the unfolding of its crown, you shall see individuals and
generations, and you shall find true the poet's word about death and life. --- ``Like the
leaves of the trees, so are the generations of men; the leaves are strewn by the wind upon the
ground, but when the forests grow green, fresh ones shoot forth from the branches at the return
of spring. So the generations of men, one is born, another vanishes''.

But it is a peculiar matter with a hope that is to be built solely on a proposition such as
this: The individuals die, but the race lives; for the proposition widens to mean: \emph{The
generations die, only the universal life lives}.

Individuals grow like grain on the field of the race; life from life, new sprouts, new stalks,
new flowers and ears: it is a fruitful field; at this sight hope shall rejoice. But the joy is
not unmixed. On the field there walks a reaper, hollow-eyed and lean; it is the Bone-man with
his scythe. Who would believe that in these dry bones there could dwell so enormous a working
power. % print: ubyre (misprint)
He mows down before his feet yellow straw and green sprouts, what he leaves is only earth and
mould. Millions upon millions of individuals, different species, different races, swarm forth
on the fruitful field; the Bone-man whets his scythe and harvests them all. Who would believe
that a skeleton could hold out like that. He gets his work done so punctually, so precisely,
that the piece
% --- p. 39 ---
\setcounter{footnote}{0}%
\opage{39}he has before him, the narrow strip of the present, is vanishingly small in comparison
with the enormous stretch he has already mown, the long, broad stubble-field of the past that
he has at his back. An ocean of gloom against a strip of light! Life sighs, memory shudders to
look at it, and hope weeps like a child.

It is not the shadow of death alone with which hope must struggle; even in fresh, present life
there is a power that is life's enemy; this power is evil. And evil, what does this one word
not enclose, what pressures of existence, what torments and sufferings, mental, bodily,
deserved, undeserved, half-deserved, do not enter as fractions under this common denominator,
evil. Not even the son of Ocean, that old Proteus, could twist himself through more windings
and assume more unlike shapes than this cunning and yet so manifest enemy of life, the dark,
enigmatic power that is the author of all evils. If man, the personal being, is to be able to
fight against evil and hope for victory, he must make an alliance, not with impersonal powers,
but with the highest personal power: he must build his hope on a \emph{governing providence}.

The need for a personal governing power, which secures the preponderance of the good in the
continued struggle and leads individuals and generations to the goal of completion, has such
deep roots in the soul that this need, even where the clear word of truth has not yet lit its
light, breaks forth with spiritual, irresistible force and gives itself form in myth. Among all
heathen, but still hoping, peoples we find a myth of providence; so among the Norsemen, so
among the Persians.

% --- p. 40 ---
\setcounter{footnote}{0}%
\opage{40}Among our Nordic fathers the struggle between the good and the evil is a constant war
between Æsir and giants. How the struggle ends is well known. After Ragnarok the evil sink down
to the serpents of Nastrond, but the good rise to transfigured life, to joy and peace in
Gimle.

Among the Persians the struggle between the good and the evil unfolds into an immense
world-drama; it is a struggle that extends equally through all the regions of nature and of
spirit, a struggle between light and darkness, Ormadz and Ahriman. By the almighty word Ormadz
created from the beginning, besides six immortal powers, a heavenly host of spirits and
Fervers. Every creature, from the animal right up to the angel, received its Ferver, its
guardian spirit. On earth was ``the high mountain'', which towered up through all the heavenly
spheres right to the first light; there Ormadz fixed his dwelling. But from the south came
Ahriman, the evil spirit, with the host of darkness, and the host of darkness was just as
numerous and mighty as that of light. Ormadz sought to preserve the peace; but Ahriman,
snorting with rage, began the strife. In the first encounter he was indeed blinded by Ormadz'
radiant splendor and by the sight of the Fervers: the word of the good overcame him; he plunged
down into the abyss. But he was not disarmed. He armed himself and gathered forces; against
every one of Ormadz's good productions he set an evil one. When he again stormed against
heaven, he was seized by such terror that in the shape of a serpent he plunged down to the
earth's center and infected everything; even fire, the purest element, he infected with smoke.
Every time Ahriman rose up against heaven, he was cast into the abyss;
% --- p. 41 ---
\setcounter{footnote}{0}%
\opage{41}Ormadz and the good spirits were victorious. But on earth Ahriman held his ground; here
he shares power with Ormadz: light and darkness, good and evil, lie in ceaseless strife. That
the struggle must also pass through human life, arousing both fear and hope, goes without
saying. The first human beings lived guiltless and pure in faith in Ormadz, until Ahriman
seduced them into drinking goat's milk and eating of forbidden fruits. It was the woman who
first sacrificed to the Divs, the evil spirits. So death came into the world with Ahriman on
account of man's first sin. But death is also a deliverance for man, for with it the struggle
against evil ceases. The souls of mortals, which were created by Ormadz from the beginning of
the world, live in heaven, whence they are compelled to descend and unite themselves with
bodies in order to complete the earthly pilgrimage. After death they are led to the bridge
Chinvad, which is broad and easy for the good, but narrow and sharp for the evil. When the good
have come over the bridge, they are led into Ormadz' kingdom, to delight and joy; but the evil
are hurled from the bridge down into the abyss to endure torments whose duration is in
proportion to their sins, but which can yet be shortened by the intercession of kinsmen and
holy men. The greater part of sinful human beings remain in the kingdom of punishment until the
end of the world, until Ahriman's power is broken and blessedness comes to embrace all.

These are only myths; but myths can, as far as the personal is concerned, contain great truths.
What chiefly matters in myth is the inward, not the outward, the life-thought itself, not its
clothing. Whether Ahriman seduced the first human beings into drinking goat's milk or sheep's
milk is
% --- p. 42 ---
\setcounter{footnote}{0}%
\opage{42}inessential; the essential thing is that an evil, tempting will seduced them. What
fruits they were tempted to eat is not the question; the main point is that they were
\emph{forbidden} fruits, that by eating of these fruits they broke the holy law, came into
contradiction with themselves, and sinned. How death in the physical sense became a
consequence of sin, of this the myth gives no explanation; the emphasis falls on the
psychological moment, that physical death, by reason of sin, appears as a hostile power that
threatens the personal life with destruction. The ambiguity in myth is that it sets up
something outward which only signifies something inward, and cannot itself give an account of
the difference between inward and outward. Mythic faith begins originally by grasping the
inward in the outward, so that a living content streams through the imagination. But
imperceptibly it happens, little by little, that the visible forms take on the stamp of
external, sensuous, independent actuality; the spirit of the myth fades away, and only the
mask, the bodily covering, remains. To be a believer at this stage is then, in the Persian
sense, to be fully and firmly convinced that Ormadz dwells on ``the high mountain'', that % print: al (misprint)
Ahriman really looked like a serpent when he plunged down from heaven, and that it is partly
the goat's milk, partly the decree of a punishing will, that has put death into the world. ---
Whoever, when consciousness, supported by the pictorial form, has reached the genuinely
dogmatic stage, does not believe everything literally, is to the religious police a suspicious
person and well on the way to unbelief.

We note that the change which mythic faith gradually undergoes is brought about by a
% --- p. 43 ---
\setcounter{footnote}{0}%
\opage{43}reflecting understanding, which throws itself one-sidedly upon the external and wants
to have all the forms of the clothing filled with solid actuality; it is this that is expressed
in the mythic events becoming certain historical facts. What now follows from this? If
reflection stops halfway, then it gets no further than changing the idea-filled mythic faith
into a certain vulgar popular superstition with an addition of unclear inwardness. With this
weaker salt, thoughts and morals and natural powers can nonetheless keep fresh for a long time.
This was the case in the ancient world with the Persians and Norsemen, whose heathenism never
got so far that it crossed the line and came under the sunburn of culture. Among the Greeks and
Romans, on the other hand, reflection got so far that it became clearly conscious of the
superstition. Once some survey had been gained of an actual connection in nature, with causes,
conditions and laws, it could not be difficult for the reflecting understanding to see that
the gods with their deeds neither were, nor had ever been, actual. Since the question concerned
something outward and ran: actual or not actual? thought, even among those who pondered the
essence of things, became more and more occupied with what is naturally necessary and turned
away from the personal.

Faith in Providence, which among the Romans was an egoistic holding fast to that decree of fate
--- sic erat in fatis --- that Rome should be mistress of the world, never came to play any
real leading role among the Greeks. The Greeks were on the whole not a hoping people; hope lay
hidden in Pandora's box and was, no doubt, really only for individuals, for private use. To
% --- p. 44 ---
\setcounter{footnote}{0}%
\opage{44}hope is romantic; it is to look at what is coming to be and to long for what is to
come. To resign oneself is classical; the classical eye does not look at what is coming to be,
but rests in the contemplation of what is, of what is complete in itself: the future is only a
continued present.

To what degree the educated consciousness, with its critical reflection, had in the second
century of our era got beyond the popular notion of a divine government of the world can be
seen, among other things, from the way in which Lucian settles accounts with the old gods. In
one of his dialogues
--- Zeus in Distress --- it begins with the father of gods and men pacing up and down his high
hall, tearing his hair and uttering cries of lament like one in despair. Hera opens the door;
Zeus wants to explain to her what misfortune has happened, but she interrupts him with the
jealous exclamation: You rogue, I know you well enough, it is another love affair again, and
goes her way. Now Hermes comes, and Zeus tells him how matters stand. You know, he says, that
some days ago we gods were invited by a rich merchant in Athens, who, when his ship was in
danger, had promised us, if it were saved, a great sacrifice and a rich banquet. He was stingy
with the fare, the miser; but since I had now once come to the city famed for its piety, I
wanted to use my incognito to look about me and learn the mood. What must I experience? In an
old colonnade, where a mass of people had flocked together, two philosophers stood disputing
about the gods and their existence. The one, a Stoic, defended them very thoroughly, and on the
first day it seemed as if the crowd would side with him. But
% --- p. 45 ---
\setcounter{footnote}{0}%
\opage{45}on the second day the opponent, who was an Epicurean, a truly abominable fellow, came
with such objections that he carried the people with him and all swore that from now on they
would no longer honor the gods, but despise them and put out the fire on their altars. These
are evil times, Hermes, it is all over with faith; what are we to resort to? Hermes advises
Zeus to call a general assembly. --- Yes, do it, Hermes, call them together, gods and
goddesses, great and small, as many as there are, say that I command it. Hermes calls out: all
gods and goddesses must appear at once; by Zeus's command a general assembly is to be held. No
one comes; all is still; how can that be, have the gods then gone deaf? You dolt, exclaims
Zeus, is that how to call the gods? You can surely understand that they pay no heed to so
prosaic a summons. You must address them in verse, a verse of Homer! Hermes begins: Hear, all
ye gods who dwell in the wide heaven\dots{} Now a rushing and roaring is heard in the air; they
come from all corners of the world and gather in the great hall beside; there arises a terrible
noise and din in there in the hall. Is it not strange, says Zeus to Hermes, I feel so odd. I
must of course give an introductory address. You know me from Homer and know that words
otherwise never fail me; but at this moment, I assure you on my honor: I am quite at a loss;
where am I to get a suitable idea? Is that all, answers Hermes; do as the orators in Athens do,
and you can never be at a loss. How then do the orators in Athens go about it? They steal most
of it from others and add a little of
% --- p. 46 ---
\setcounter{footnote}{0}%
\opage{46}their own. The art is: not to take too much from any single one, but to borrow a little
from the one and a little from the other: then it is not plagiarism. --- Excellent, Hermes,
that I will do. Go in meanwhile and bring them to order. --- Hermes goes, but cannot get the
better of the confusion. All are pushing forward; everyone wants a seat among the highest.
Hermes goes back to Zeus. ``Do not disturb me, I am now in the middle of my speech''. But it is
not possible to keep them in hand; I must have a rule: our Olympian ordinance of rank is sorely
in need of revision. Yes, says Zeus, if justice were to take its course, they would of course
have to rank according to the masters who have shaped their images, the Phidias gods first and
so on downward. But the age is materialistic; we shall have to go somewhat by the material. Let
the gods and goddesses who are of gold, silver and ivory sit uppermost; you understand me.
Hermes goes with this message; but a new quarrel arises. Apollo speaks of his ``silver bow'';
Aphrodite claims that she is as good as of gold, since Homer has called her ``the golden'', to
which another remarks that her gilding is a peculiar matter. So it goes on, until the Rhodian
Colossus steps up. He is, as he says, neither of silver nor of gold; but he is fully conscious
of his superior power; he has it in mass. He bangs on the table and asserts the right of mass.
This impresses the assembly; it comes to order. When all have been seated, Zeus enters, clears
his throat and gives the introductory speech. After describing the situation to them and
showing them the misfortune in all its magnitude, he ends by giving the floor to whoever might
first request it. General silence. Repeated invitation;
% --- p. 47 ---
\setcounter{footnote}{0}%
\opage{47}suspense and silence. At last someone rises; it is Momus. I am glad, says Zeus, that you,
Momus, are in this serious moment so well disposed that you can forget your mockery and give us
good counsel. You have the floor, let us hear your counsel. Good counsel, says Momus, I have
none; but if you gods will hear a word of comfort, then it is this, that if men now turn their
backs on us, then we must all surely know in ourselves that we have honestly deserved it. When
have we gods ever in earnest concerned ourselves with the affairs of men? Bethink yourself,
Zeus, what masses of violence and injustice lie heaped up on earth, steadily growing all the
time you have held the scepter. Where were you when freedom in Hellas, fighting, dying, drew its
last sigh, when the Achaeans had to send the thousand to Rome? What were you doing the day
Mummius took Corinth? When you were not sitting nodding on one of Olympus's ridges, you were
running mad after the nymphs and the pretty girls. Hear! hear! --- from Hera. Yes, Hera, you
too the human race has no particular need to praise to the skies; your mind is like the dark
air, you have always followed your hatred and your whims. --- In this spirit, if not with the
same words, Momus settles accounts with the myth of providence, and consoles the gods with
this, that even if there are individual thinkers who get so far as to see through their game,
there will still always be an ignorant mob that will surely go on worshiping them and hoping
for their help and assistance.

Against the Christians Lucian turns with, if possible, still greater scorn. These wretched
people imagine that the Deity concerns itself with them; they croak like frogs in a swamp and
boast of a providence
% --- p. 48 ---
\setcounter{footnote}{0}%
\opage{48}that counts the hairs on their heads, so sure of themselves and yet so credulous that
they let themselves be deceived by every juggler.

The old, hopeless world with all its knowledge, all its wit, went step by step to the grave.
Great in its strength, wonderful in its patience, invincible despite indescribable sufferings,
the new faith with the new hope raised its banner high; and the banner became a sign of power,
an emblem of eternal life, a light of blessedness in the night of death, a shield against all
dangers.

It is now almost sixteen hundred years since the powers of the world, in civilization and
barbarism, began to bow before the new sign; and yet the old strife is to repeat itself once
more, though under other forms. Modern civilization, which has turned more and more toward the
impersonal essence of nature and chosen Greece and Rome as its ideals, has, as we are assured,
become convinced that the hope of blessedness is illusion and dream, and the thought of an
almighty, free, all-governing will not a thought but a fancy. The relative is a moment in the
absolute, the corruptible inseparable from the incorruptible: evil is just as necessary as
good. Without vices no virtue, without shadows no light, without death no life; no truth
without lies, no joy without sorrow, no enjoyment without suffering, no force without
resistance, no symphony without dissonances, no drama without a villain: that is, in brief, the
new wisdom. From this it follows with merciless logic that even if the good must be posited as
the goal, evil is nonetheless necessary for the sake of the good. The dynamometer is the same
for both: if the good grows, then evil grows too. But from this
% --- p. 49 ---
\setcounter{footnote}{0}%
\opage{49}it then also follows --- with merciless logic, that every hope of a better future for
individuals as for the race is vain folly. If there is not a veiled contradiction, a secret
flaw in this logic, then it follows incontrovertibly that the new world, when it again reaches
the height of civilization, must go to its grave as hopeless as the old. This hopelessness has,
moreover, in the very latest time found its scientific expression in so-called nihilism. The
riddle is that all is from nothing, and the riddle's solution that all returns to nothing. What
then has the race to hope for? Nothing. The wisdom is that the race, when it has seen
everything, tried everything, seen through everything, must at last be done with everything and
discover the goal: nothing. Like an old man, full of days, weary and tired, the race that has
lived out its life will take leave of life with a melancholy smile and be forever reconciled
with death.

Nihilism is dark and melancholy; but in its dark melancholy it has something most serious and
purifying. It is finitude, temporality, the many-colored interests of the motley world and the
vanity of its great works, on which it sits in judgment. The finite, held fast by itself, is
illusion, and the ground of illusion is nothing. As the whole of existence thus comes into the
crucible of all-testing criticism and is melted down, nothing is lost but the slag of finitude.
The result is far from being as hopeless as it might seem at a first glance. The great nothing
that swallows everything cannot possibly be empty; it is the ground of a spiritual possibility,
into which everything sinks; the whole is still there, but transformed in such a way that it
must anew be grounded and determined by spirit.

% --- p. 50 ---
\setcounter{footnote}{0}%
\opage{50}Nihilism is not nearly so dangerous to hope as the otherwise so favored expectation of
an impending great and blissful future. It is a view that builds securely on the ground of
finitude. Now we must certainly grant that finite ends and temporal interests have their
reality, and that it would be sad if these interests, with all the energy they call for, were
to be weakened and paralyzed by a melancholy speculation. It is not finitude, but its ground and
grounding, of which we speak. It is the incontrovertible truth, which must be impressed again
and again, that temporal life never reaches and never can reach a full, harmonious conclusion
when its final goal is set in a temporal time. The expectation is that civilization will one
day, perhaps by a coup de main, perform the great miracle that all strife and discord cease,
that all receive an equal share in the goods of life and therefore live in peace and good
understanding. That exceeds the powers of civilization: a hopeless hope! Sooner shall the ocean
of the world fall asleep into eternal rest, all waves lie down, all rushing currents stand
still, than temporality, this corruptible being whose innermost is unrest, strife and discord,
shall come to peace, to harmony with itself, and settle into a state of equilibrium. No,
struggle and struggle and again struggle, that is, in the name of eternity, the watchword of
all temporality. In a temporal sense it can be said with reason that evil must serve the good.
The stir of society, the wild cries, the dull din, the hurrying footsteps, will surely disturb
the sleep of effeminacy, the beautiful dream of ``the civilized'' of an earthly bliss.

There must --- for eternity's sake --- remain in the minds of men a secret dread, an open
wound, a
% --- p. 51 ---
\setcounter{footnote}{0}%
\opage{51}rift; for only through this rift can the sigh press forth, the sigh of longing that
divines something incorruptible. There must be cracks and fissures in the masonry of
temporality; only through these cracks and fissures does a gleam of light fall in from the
eternal.

Could a worldly prudence succeed --- if only for a little while --- in hushing all complaints,
stopping the mouth of all discontent, sating even greed itself, and getting the rifts and
fissures of earthly well-being duly puttied up, what would then happen? Then the spiritual tone
would vanish; the last greeting from higher regions, the last faint echo of the harmonies of
those spheres, would die away. It would go with an existence obdurate in finitude, petrified in
temporality, as according to the legend it went with that marvelous statue --- the statue of
Memnon. Every morning, when the sun rose, the statue gave forth a tone; it sang in the stone.
There was something spirit-like in it. People listened and perceived the divine: it sounded
like a greeting from above. The secret was that there was a crack in the stone, a fissure in the
statue. The fissure was discovered, but the secret was not discovered. It was thought that when
the statue was repaired, it would really begin to sing and give forth tones in full harmony. How
disappointed they were; for what happened now? ``Alas, the statue fell silent, and the sun
poured its silent gold into the basin of the desert sand. And never more that clear greeting
sounded, from the same hour that the crack was healed''.

% --- p. 52 ---
\setcounter{footnote}{0}%
\lecture{Fourth Lecture.}\opage{52}It is a not uncommon opinion that positive, revealed religion contains a number of articles
which do not really concern faith, and which one may therefore leave aside; but, it is said, there is one article
which for the sake of life and hope must absolutely be preserved, and that is faith in Providence. The doubting
understanding may take many articles away from a full ecclesiastical confession of faith, but the one concerning
the fatherly providence must remain standing; it is the most important article of all. There are those who add one
more article which is also to remain standing, namely faith in the immortality of the soul; and there are those
who, to be on the safe side, add yet another article, namely faith in a righteous yet merciful Judge, who will at
last settle the account of our life and pay every man according to his works. When these three articles are in good
order, one supposes one has all that one needs of positive religion. Now how is this to be understood?

It must be granted that in the presentation of phenomena and facts through which the content of revelation meets us
from without there may
% --- p. 53 ---
\setcounter{footnote}{0}%
\opage{53}occur much that lies remote from us, because it is in historical form and belongs to a remote past, much
that at first glance repels us, because our view of the world, the picture of the reality of heaven and earth to
which the culture and thought of the present have naturally led us, is so widely different from that of antiquity.
It must be granted that the ecclesiastical, indeed even the biblical tradition, in its several links and joints, is
far from furnishing an outwardly secure foundation, a dogmatic framework of beams, to which one can fasten faith and
hope as with spikes and nails; all this may be granted, and still the view of positive, revealed religion indicated
above is burdened with great errors.

The very circumstance that there is talk here of articles, one, two, three articles, and so of a parceling out, must
surely give one pause. If revelation contains truth, then the truth does not let itself be parceled out. Revealed
religion, in its essential content, is either false from first to last or true from first to last; there is one
spirit, one faith, one hope in all its proclamations, one stream of life runs through the whole. He who has in truth
appropriated the revealed word about Providence, the fatherly providence, without whose will not a sparrow falls to
the ground, has also appropriated the word about sin and grace, about fear and hope, about eternal life,
righteousness and judgment. He who has not appropriated the latter is in error when he imagines that he has
nonetheless understood and appropriated the former. Positive, revealed religion is from first to last a revelation
of providence: all is gathered in the one faith, faith in Providence. Positive, revealed
% --- p. 54 ---
\setcounter{footnote}{0}%
\opage{54}religion is from first to last a revelation of righteousness and of love: all is gathered in the one faith,
faith in a final personal glorification of all that is good and a final personal condemnation of all that is evil.

How do we convince ourselves of the truth of this? What gives us certainty? The source of life, life itself! We
first live as we learn to live. The measure of faith is experience; the measure of hope is experience, and --- it
cannot be repeated too often --- the measure of the appropriation and understanding of revealed religion in its
wholeness is, from the beginning to the end, experience. We hold fast to the outward only in order through it to
appropriate the inward; but the appropriation is in experience. Whatever in what has been handed down historically is
of such a nature that it cannot in any way, either immediately or mediately, become the object of inner, personal
experience, is no concern of religion, can be neither an object of faith nor a ground of hope.

By applying the measure taken from experience itself, or at any rate from its possibility, we become able to
characterize the various misunderstandings of the doctrine of providence and, from the standpoint of religion, to
contribute to an elucidation of the living Providence's self-revelation to individuals and generations.

Among the misunderstandings we then meet, in the first place, the superficial or vulgar faith in Providence. It is
composed of two elements: the notion of a general world order, on the whole unchangeable, and the notion of a
``finger of God'' which on given occasions intervenes where it matters. That such
% --- p. 55 ---
\setcounter{footnote}{0}%
\opage{55}two notions cannot straightforwardly be reconciled in one consciousness is understandable. In experience,
accordingly, they are mixed in the most various ways.

If the notion of an unchangeable world order has the upper hand, then the doctrine of providence sinks down to a mere
human invention. If we knew all the circumstances, we should be able to calculate the occurrence of events as an
astronomer calculates a solar eclipse or some other constellation. Now we do not know the infinite concatenation of
causes, conditions and effects. We come upon chances, incidents, remarkable coincidences. If these coincidences of
events have great, beneficial effects either for ourselves or for the whole, then we are moved and, in the moment of
being moved, believing: we refer them to Providence; if they seem to us insignificant, then they are mere incidents;
if they weigh upon us as adversity, then it is fate; if we come to calm, cool reflection, then it is the natural
course of things. Our faith in Providence and our hope are grounded on our ignorance.

If the notion of an almighty, arbitrary will has the preponderance, then man places himself in a kind of fairy-tale
relation to Providence, ventures everything, even the most foolhardy, on Providence's account. Just as a frivolous and
heedless son who has a rich and powerful father permits himself all manner of excesses, trusting that the father, who
is somewhat strict, to be sure, but at heart a good and loving father, will surely get him out of the difficulty and
come to his aid when things get tight, so men can become fantastically believing and build a filial, fantastic hope
on the fatherly finger of Providence. It is incredible how easily they become believing and
% --- p. 56 ---
\setcounter{footnote}{0}%
\opage{56}in faith hoping, when it agrees with their inclinations. A spendthrift can live in style on what he has, if
he does not trouble himself at all about the future but lets Providence take care of the morrow. A whole nation can
save millions when it has ``the right, living faith'' and puts its trust in the finger of Providence. Its system of
defense, for example, will not cost it much; it is in no hurry at all to put it in order. Should a lurking enemy one
day be tempted to want to lay his fingers on so heartily believing a nation, then let him beware lest Providence lay
its fingers on him: such a faith with such a hope is a great saver.

A couple of traits from actual experience. There was once a king who lived through the heartache of having to make
peace with an overpowering enemy, after having ceded to him about half of his kingdom. It was a hard fate; but it
could of course also be understood as a dispensation of Providence, and was at any rate understood so officially.
Divine service with prayers of thanksgiving was ordered throughout the land. The king with his court was in church.
Behind the king's chair sat a courtier, mumbling during the sermon. The king turned round and asked: what are you
mumbling about? I am talking with our Lord, Y. M. What are you saying to our Lord? I am thanking him for the peace,
Y. M. And what does he answer you? Don't mention it! Here we have an example which shows that the notion of the will
of Providence has gone under in the notion of the laws of the world and the natural course of things. --- Another
trait falls under the pusillanimous. There was a decent woman whose husband was a virtuoso. This same woman at certain
times, owing to circumstances, came into a peculiar, quite
% --- p. 57 ---
\setcounter{footnote}{0}%
\opage{57}singular relation to Providence. The virtuoso, that is, was in the habit of giving once a year a large,
extraordinary evening entertainment. On this occasion the wife went to church, usually twice. The first time, before
the evening entertainment, she prayed to Providence, and her prayer was a prayer of petition; she prayed from the
heart, fervently and sincerely, that it might go well, she prayed for a large attendance, luck and success on the
violin, etc. The second time, after the evening entertainment, she prayed to Providence, and her prayer was a prayer
of thanksgiving, that is to say, when it had gone well; for when it had not gone really well, she did not give thanks
at all. Here is an example which shows that the notion of the natural course of things can, when it matters, be
displaced by the notion of a supernaturally intervening will. The pusillanimity does not lie in the fact that the
anxious woman prayed for something temporal; for if it were not permitted to pray for something temporal, then it
would not be permitted to pray for our daily bread either; no, the pettiness lies in the selfishness, that she forgot
to give thanks when it had not gone as she wished.

The next misunderstanding stands on a higher level. Here one builds, quite rightly, on the presupposition that the will
of Providence does not go around the laws but through the laws, and is master of all, the least with the greatest.
But then the demand is also made that the assistance and help of Providence shall make itself clearly known in this
actual, present life. The notion of a beyond is either quite forgotten or at any rate pushed back; faith in Providence
is taken as an article by itself, separated from faith in immortality. From this point of view
% --- p. 58 ---
\setcounter{footnote}{0}%
\opage{58}what Leibnitz says holds good: ``it is no small matter to be content with God and with the universe.''

It comes far easier to the selfish in man, with all possible zeal for righteousness, love and mercy, to appeal to
nature and resign himself to necessity than to hold fast to faith in an all-governing will and to hope for the
assistance of Providence. No one who is in danger will be angry with the powers of nature, even though they threaten
him with destruction; they cannot do otherwise, they are lawfully excused. No one who is in distress at sea becomes
embittered either at the storm or at the waves or at the rocks; the will-less beings bear, after all, no guilt for the
devastations they cause. But if there were one out there with them of whom all knew that he could help if only he
would, and yet did not help, he would become the object of a furious bitterness. They call on him, they pray, they
implore him, they cry out to him, all in vain! He stands silent, unmoved, cold as the wave, stone-deaf as the storm,
hard as the rock, --- a heartless monster! The struggling mock him, the dying curse him, and those who survive the
danger declare that it was no real person but a phantom, a vision in the mist, an offspring of an excited fancy. So it
has gone with many people in regard to Providence. They began by believing, they came into danger, they hoped, they
cried out for help; but the Silent One gave no answer, the help did not come. Either they then sank into the deep with
a curse on their lips, or a chance helped them: they ran from the danger, thanked the accident and mocked the
impotence of Providence. It was a fault in the experience; there was
something they never got to see in the danger, a sight that faith
% --- p. 59 ---
\setcounter{footnote}{0}%
\opage{59}alone can show to hope. ``He stood, the Silent One, on the ship out there, in the midst of the howling storm,
and saved them all, --- not \emph{from} destruction, no, \emph{through} destruction.'' It was faith that discovered it,
for the eye of faith is sharp, it sees through mists. Faith saw at the rim of the horizon the new coast, the quiet
harbor. That sunny land rose into view; the shadows of the groves, the glinting silver of the springs, even the
flowery bloom on the green meadows faith could see. And faith became glad and whispered to hope: ``see how he, the
Silent One, he whom they curse, stands and saves them all, each and every one; no one is forgotten, neither small nor
great, there is a way, there is room, there is a dwelling for them all. So we two, you and I, will not curse him; no,
you shall praise him aloud, and I will fall down and worship him.''

The third form of the misunderstanding, a modification of the preceding one, is characterized especially by selfish
impatience. The one who hopes does not give up his faith in Providence; but there must be right and justice in the
world: what becomes of righteousness, where is the retribution? It goes so slowly. Wickedness gloats, the good is
oppressed, the strong man boasts of the strength of his arm: it looks at times as though Providence were asleep.
Awake, how long shall it last? That is the cry and complaint of impatience. The two articles, faith in a governing
Providence and faith in a righteous Judge, are here put together to the exclusion of the deeper thing that lies
between them. It is constantly forgotten that the final settlement is not here but yonder; it is forgotten that
temporality is not merely an educational institution but above all an arena of freedom; it is
% --- p. 60 ---
\setcounter{footnote}{0}%
\opage{60}forgotten that the last act of the great play falls not within but, in an essential sense, outside the
temporal boundary of life.

In the theater, when a great drama was performed that really put the minds and thoughts of the spectators in
suspense, it is said to have happened at times in older days, when naivety was greater than it is nowadays, that an
unprepared but thoroughly honest and guileless person disturbed the performance by his impatience. In such a play
there are naturally evil passions, falseness, slander, misunderstandings. There are scenes, clashes, where a single
word would be enough to clear the whole thing up, if only it were spoken. The honest man cannot bear to watch it. He
wants to betray the secret; he cries out aloud and wants to come to the aid of the misled. They hush him and get him
quieted; but it does not last long, for in the fourth act it gets worse and worse. It is he, ``the villain, the
red-haired scoundrel with the bushy brows and the black eyes,'' whom the honest man cannot for any price abide; ``for
it is he who goes about and causes all the misfortunes.'' The honest man is in revolt, his impatience tolerates no
objection; he shames the lukewarm spectators for sitting there so indifferent; he calls for the police and wants to go
straight up on the stage and make things plain before it is too late; he forgets what stands on the playbill, that the
play has --- a fifth act. That is an image, if an imperfect image, of human impatience toward the governing
Providence. We trust the poet when the play takes its course and holds our expectation in suspense; but toward the
great world-drama,
% --- p. 61 ---
\setcounter{footnote}{0}%
\opage{61}in which we ourselves have a small part, our patience gives way: we are indignant at our fellow players and
take offense at the invisible poet. What scenes when a Roman triumphator makes his entry into ``the eternal city,''
what an image of righteousness on earth! We can scarcely bear to read of others' misfortunes in history; if it
concerns ourselves, we easily despair of the righteousness of Providence. We forget that the play has a fifth act,
which is played on another stage. --- These misunderstandings must be removed if it is to be possible to appropriate
the word of revelation about a living Providence.

The essence of Providence is will, almighty will. The great drama in which every man plays a part is designed to
develop the human will and bring it, through strife and errors, into agreement with the divine will. What, then, does
this divine will will, what is its real plan with man? If the plan lay in the unconditionally closed necessity of a
system of reason, it would have to be discoverable, in its main features at least, by thorough scientific thinking.
Here no special revelation would be needed. But since its innermost principle is will, freedom, an unconditioned
freedom which makes the system of reason its foundation and posits everything in nature as means to its end, it is
impossible that this end, the end of the unconditioned will, can become revealed to man otherwise than by the will
itself revealing it.

This revelation has two stages. The first stage is marked by the laws, the great laws for obedience and freedom. The
laws concern the general; yet, for the sake of life and actuality, the particular is also entered into with definite
commandments, threats and promises,
% --- p. 62 ---
\setcounter{footnote}{0}%
\opage{62}punishment upon the wicked, reward and blessing for the good. The retribution proceeds from the Eternal, but
yet is couched in terms of the temporal; so the thought of retribution awakens a painful strife between faith and
doubt, between hope and disappointment. Hope is disappointed and yet holds out in spite of all disappointments; it is a
continued crisis, but a crisis of faith and obedience, a crisis during which the tension grows and hope is purified.
The second stage is the fulfillment of hope, the coming of the Expected One, the personal self-revelation of
Providence. What this manifestation means, in what relation it stands to the goal of life, the salvation of
individuals and the development of the race, can be seen, among other things, from the way in which the Expected One
was received. ``He came unto his own, and his own received him not''; these words remind us of a strife that began at
the turning point of the ages and will continue until the end of the world.

``Is that to be Providence revealed? Is this man God himself, does God look like that? Is that his thought, his way,
his means and end? A reasonable person will find that hard to get into his head. To be sure, the race might sorely
need a providence that would help the oppressed against their oppressors and take charge of affairs that have fallen
into confusion; but to offer oneself in such a form leads to nothing. If Providence is to break through and take
effect, it must come in a far grander form, employ other means and in earnest seize the reins of world government.''
We cannot by any objective general consideration or scientific demonstration whatever convince men that the Expected
One really is Providence revealed; but there is nonetheless something of which all who think can
% --- p. 63 ---
\setcounter{footnote}{0}%
\opage{63}convince themselves, and it is this: either the self-revealing Providence must go the way he went, or there
is no Providence, the race is hopeless. That can be proved to the letter.

The riddle lies in the relation between the divine and the human, the almighty and the powerless will. This riddle is
not solved by an admonition such as that man is to obey God's will; for when obedience is placed in blind submission to
outward commandments, freedom is lost: there is no essential self-determination, no content of experience, no
development. The riddle cannot be solved by referring man to his nature and natural fundamental will; for if the
starting point is sought in nature, man must follow his instinct and is given over to obscure impulses; if man will
pull himself together, work up the natural impulses and use his finite will as it is, he must rely on himself, and
what then does Providence mean? Nor can the riddle be solved by pointing to the idea of the good, for the idea of the
good in ``immanent existence'' is a scientific myth. The idea, it is said, determines the will, the will is absolutely
dependent on the idea, i.e. the will determines the idea, the idea is absolutely dependent on the will. In separation
from the idea the will is arbitrariness; in separation from the will the idea is necessity. Freedom is the feigned
movement whereby necessity is absorbed into arbitrariness, so that arbitrariness may in turn sink down into necessity.
Immanent freedom is a self-contradiction without hope; here there is no Providence.

Let us consider the Expected One. What is he? An actual human being in flesh and blood. Whence comes
% --- p. 64 ---
\setcounter{footnote}{0}%
\opage{64}he? From the Father's bosom, from the holy deep where the Creator mirrors his image in the race, where God's
being is one with and yet distinct from man's being, where the will itself is nature, and nature holy will. So he
stands there, on the stream of the ages, the revealer of the Almighty and yet powerless as a man. It is
incomprehensible, they say. Right! It is nothing to offer to science, for it is, to put it plainly, a wonder. But this
wonder is precisely freedom, for freedom itself is a wonder. Should we now reject the wonder because it cannot be
grounded in objective knowledge, then we must at the same time reject freedom, providence and hope, for freedom,
providence and hope are not scientific quantities.

The wonder of freedom presupposes a peculiar relation of will between omnipotence and impotence. The divine will is
free, for it is almighty. Freedom is inseparable from power, unconditioned freedom inseparable from omnipotence. The
human will is unfree, for it is powerless: unfreedom is inseparable from impotence. In the manifestation of Providence,
then, we are to see the wonder that the almighty will makes itself powerless, that the powerless will can become
partaker of the power. Impotence is clothed with omnipotence insofar as man's will becomes the accomplisher of the
divine: therein consists, in the absolute sense, the \emph{freedom} of the human will.

To will what Providence wills is to will what the Almighty wills, is to have an irresistible will. The exemplar,
Providence revealed, has made plain in great facts what this means. He stands at the grave of Lazarus, and how?
\emph{Weeping}, \emph{praying}, \emph{commanding}. To weep and to pray is human,
% --- p. 65 ---
\setcounter{footnote}{0}%
\opage{65}it points to impotence; but the voice that commands: Lazarus, come forth! is divine and points to
omnipotence.

Now if we fix our eyes wrongly on the outward, then such a miracle easily becomes an offense. ``If Providence was
revealed in order to raise the dead, why then did it stop at a few single, highly doubtful attempts and go no
further?'' We will not here review the host of objections with which scientific criticism stands ready to attack; we
are done with them all, in granting that they may well have their validity from the standpoint of an objective
knowledge; our whole refutation is compressed in the declaration: \emph{the standpoint of objective knowledge is
false}. The question on whose answer everything essentially depends, with this as with every other miracle, is
directed not immediately at its actuality, but first and last at its \emph{possibility}. It is a boundary of
possibility that categorically divides faith from unbelief. If a miracle is \emph{impossible}, then there is no free,
almighty will, no unconditioned principle of freedom, no living Providence, no hope: then unbelief is right. If the
miracle is \emph{possible}, then there goes through the outward a way to the inward, and then an experience is
possible too.

Held fast as a historical fact, the raising of Lazarus cannot be experienced by us; the only thing we can experience is
an endless, confused strife about the actuality of the fact. If, on the other hand, we look to the inward, then
something else shows itself here. The word of omnipotence that sounded at the grave of Lazarus is a manifestation of
God's secret, that the innermost principle of all existence
% --- p. 66 ---
\setcounter{footnote}{0}%
\opage{66}is freedom, that the almighty will is, within the coherence of the laws, the Lord of the laws, so that the
highest that man's will can attain is to become the organ of the almighty will. It is a great, divine contribution to
the clarification of the whole thought of providence that the miracle offers us. What weighs down, what cows the
upward-striving will, awakens despondency and casts down hope, is, from the standpoint of finitude, always this: how
will freedom be possible in this physically bound existence? Provide possibility! The more the will exerts itself and
its powers are spent on attaining a higher good, the more obstacles pile up on the way; it may look as though blind
laws barred the possibilities and the unfree powers conspired against freedom. Now what does the word of Providence at
the grave of Lazarus say? Fear not! If you will only what Providence wills, then there is possibility enough,
possibility for faith and hope.

Can, then, it is asked, the human will, by being the organ of the will of Providence, itself become almighty? By no
means; every selfish thought must go. If there is but a remnant of selfishness left in him in whose hand Providence has
laid the execution of a great plan, the self-willed will, cunningly or openly, will seek to seize the power for itself,
use it in its own name, for its own glorification, i.e. misuse it, and so at last collapse into impotence. Religion
uses for this the expression that such a one will not give God the glory. The expression is not seldom misunderstood,
as though Providence were chiefly concerned with outward marks of honor. This misunderstanding furnishes an excellent
cloak for egoism. The Roman triumphator went to the Capitol to give thanks to the gods;
% --- p. 67 ---
\setcounter{footnote}{0}%
\opage{67}the Christ-confessing triumphator goes to church, sings the \textit{Te Deum} and hangs up the trophies of
plunder and violence to the glory of God: the difference is not great. The only condition on which the powerless will
can be in true unity with the almighty and work as an instrument of omnipotence is obedience, unconditioned obedience.
Unconditioned obedience means for man devotion, unlimited self-sacrifice in the service of Providence.

Why did the Savior weep at the grave of Lazarus? We answer: because the thought of the Savior, the great thought of
providence, in which omnipotence is united with impotence and demands of man his sacrifice of obedience in humiliation
and suffering, stood living before him. If help could be given by a word of omnipotence, then he knew that the Father
in heaven would hear his prayer. And the prayer was heard, and he stood before the eyes of the people as the Almighty;
and that meant for him that now the measure was full, now his enemies swore his destruction, now the manifestation of
impotence was to begin, now he was to drain the bitter cup in humiliation. To raise the dead man was not the help; the
miracle was a consecration to suffering, a sign of omnipotence that all that he suffered he suffered in voluntary
obedience: the help was that he himself went into death.

Omnipotence and impotence, freedom and obedience are united in love: the manifestation of Providence is the % print: Aaabenbarelse (misprint)
manifestation of the divine love. In the spirit of love, impotence itself becomes almighty.

Upon all the love that stems from spirit a divine stamp is set, namely --- superiority. Even when the powers of the
spirit stand in the service of self-love, they can still, especially when the struggle has been for something general
and ideal, firmly establish a higher
% --- p. 68 ---
\setcounter{footnote}{0}%
\opage{68}self-esteem, a natural pride, which does not deny itself even though the hero, overwhelmed by enemies, is
maltreated and martyred to death; but a superiority of this kind gets no further than resigning, keeping silent before
brute force and despising the wickedness of men. He who was formerly mighty is now powerless, and the eloquent silence
of resignation says that he is conscious of his impotence. It is otherwise when the powers of the spirit stand, not in
the service of self-love, but in the service of universal love. The King of Truth with the crown of thorns reveals the
absolute superiority of the eternal love. He has not merely the fortitude to keep silent and suffer; no, in the sense
of inwardness he has the powers of an omnipotence to spare. In the bitter moment of insults, pains, the agonies of
death, he has the infinite superiority to be able to raise his voice and call down forgiveness upon his enemies. That
one who is powerless can, in voluntary sacrifice, be so almighty, that is a manifestation of Providence, that is the
incomprehensibly great wonder of universal love. He who will weigh the wonder in the two words, the word of omnipotence
at the grave of Lazarus and the word of forgiveness from the cross, shall surely find that the latter, when it is
weighed on the scale of experience, is equal to the former.

Providence revealed is the exemplar of the race, its Savior and Leader. The will of Providence is not only divine but
also human. Providence has not remained seated in heaven and prescribed for poor mortals laws that are irreconcilable
with flesh and blood, laws of the spirit so cold that they kill natural individuality, laws of omnipotence before
which the powerless must shudder; Providence has itself put on flesh and blood, has revealed its omnipotence in the
form of impotence, has fulfilled
% --- p. 69 ---
\setcounter{footnote}{0}%
\opage{69}its own laws and lived a human life on earth. It is not an arbitrary bursting of the coherence of things, not
a breach of the laws of nature, but freedom, freedom itself become man, that is the miracle. The exemplar is not only
outward in history but also, though obscured, inward in the race. The living Providence speaks to men not as to
individuals of a lower order but as to brothers and sisters; and where the word, the word of freedom, of obedience, of
love, finds entrance, the image rises in clear lineaments. This being is not foreign to our nature, it is the holy
ground of our race; it is in the souls within, if only it may come forth to inner experience, to the clarification of
the goal of life, to the fulfillment of the thought of Providence, to the firm grounding of faith and hope.

% --- p. 70 ---
\setcounter{footnote}{0}%
\lecture{Fifth Lecture.}\opage{70}If human life is to be concluded within the bounds of temporality,
then faith in Providence is an illusion, and the race essentially hopeless. It has often been said,
and it is still said, that the hope of immortality is only an offspring of the individual's childish
fear of death, a selfish craving to preserve its dear, vain I with all its limitations into infinity.
``It is self-love to lay such weight on conserving one's individual self. Instead of clinging
to an empty hope of a personal immortality, we ought to endeavor to lay aside our egoism, forget
ourselves, work for the great whole, seek rest at the breast of nature, breathe out into the universal
and see to it that we become moments in the Absolute.'' Hence: \emph{the impersonal
is higher than the personal}.

If this is now true, then he who calls himself the Saviour of the world, the exemplary personality
on whom believers personally build their hope, in whom they see Providence revealed, cannot possibly
be the truth. What is he then? A most wonderful, profound enthusiast, certainly a religious genius
on an extraordinary scale, but a mortal, limited individuality,
% --- p. 71 ---
\setcounter{footnote}{0}%
\opage{71}caught and captive in much egoism. The thought of life which the Exemplar came into the world
to reveal can be expressed in the single proposition: \emph{the personal is, in
heaven and on earth, higher, infinitely higher
than the impersonal}.

If it is egoism to place one's life in the personal and lay weight on oneself, what mortal on earth
has then drawn such attention to his personal self as he who first and last and always, with the
emphasis of divine authority, says \emph{I}. No prophet has gone further than to begin with
a: \emph{Jehovah says}; but this one, the Prophet of the prophets, begins from himself with a: I say.
``Verily, verily, I say unto you.'' ``Before Abraham was, I am.'' ``I am the truth.'' ``Come unto me,
believe in me, abide in me.'' ``I am the vine, ye are the branches.'' ``Learn of me!'' ``I came down
from heaven, not to do mine own will, but the will of him that sent me.'' ``I have life, not that
I should preserve it, but that I should sacrifice it; as a corn of wheat it shall fall into the ground,
I shall give it up for the salvation of many.'' If this is egoism, then we must admit that it is an
egoism of a peculiar kind.

There is an easily recognizable difference between the egoistic and the personal. The egoist holds
on to his I in a finite sense; he renounces no right and brings no sacrifice unless he sees that he
has advantage from it. He can venture everything, life too, but always under the presupposition:
\emph{to be or not to be}. Precisely because he never carries it further than to conceive his I as
finite, he can even with a gambler's boldness put it at stake; if it is lost, what more is
% --- p. 72 ---
\setcounter{footnote}{0}%
\opage{72}there to it? What is finite must have an end. Personality lays infinite weight on itself;
for the self knows itself as eternal; the personal assurance of the self's eternity is expressed in
deed by man's winning himself by sacrificing himself: this circular movement through abasement to
exaltation, through death to resurrection, through voluntary self-sacrifice to infinite
self-affirmation, this ring of eternity in which personality circles about its original center, is
incomprehensible to egoism. There has often been complaint that objective science, namely philosophy,
is so inclined to set the personal lower than the impersonal, that it will not acknowledge God as a
personal being. These complaints stem from misunderstandings of the nature of science as well as of
personality. There is no lack in the history of philosophy of well-meaning % print: Philosphiens (misprint)
attempts to get the doctrine of God's personality scientifically proven in objective knowledge. The
concept of personality, it is said, is the highest, richest and fullest of all concepts; consequently
no system of concepts can be completed and concluded except in and by and through the concept of a
personal God. We need not immerse ourselves in intricate analyses of concepts in order to come to the
recognition that a great error hides herein.

To a full concept of God's personality there belongs, as positive, revealed religion emphatically
impresses upon us, something other and more than universally necessary determinations of thought of
a self-being infinite in power and wisdom. Even if the science of ideas is able to clarify what spirit
is in its universal concept; even if the doctrine of spirit, as we assume, will bring it so far that
even the objective men of experience at last see that
% --- p. 73 ---
\setcounter{footnote}{0}%
\opage{73}the being which bears nature and grounds its formative drive and the wisdom in its
unconscious productions must be a seeing, knowing, almighty willing spirit: yet with such a concept of
spirit we are still infinitely far from the recognition of what religion understands by God's
all-fatherly personality.

Religion speaks of God's wrath, mercy, grace, and so forth; those who assume that such divine
emotions can be treated and grounded in objective knowledge must be ignorant of what science is. It is
easily said, superficially seen also easily understood, that the divine being cannot be set in
passionate turmoil in the same way as a human being. In the name of reason, men have endeavored to
get the Jehovah of the Old Testament freed from the terrible passions ``that disfigure his countenance
and make him so vehement that his wrath becomes a consuming fire''; these are, it is said, figurative
expressions, human representations. ``God is far too reasonable to be able to become angry. ``If one
will not remain standing halfway, then one must necessarily concede that God is also too reasonable to
let himself be moved by self-inflicted misery, far too exalted to engage with individuals one by one;
they are, after all, in a reasonable world; let them sin, rage, repent, regret, whimper, pray and bear
the consequences of their actions, until they learn to come to terms with the eternal, unchangeable
wisdom. So reasonable, so exalted, so unchangeable, so completely freed from all affects and emotions
must the true God be, if he is otherwise to enjoy the honor of being introduced into science.

% --- p. 74 ---
\setcounter{footnote}{0}%
\opage{74}Conceived in the form of knowledge, the Godhead is a being without feeling, without mind,
without heart. That such a being is impersonal is seen also from this, that with all its knowledge,
power and willing it is to such a degree occupied with the whole, the necessary, the universal, that it
cannot possibly enter into any particular personal relation to individual personalities. This
conception of the divine being is now, according to positive, revealed religion, a soul-corrupting
fundamental error. To remove this fundamental error and reveal God's personality, the Saviour of the
world has stepped forth as the Only-begotten of the Father; it is in this sense that he lays weight on
his I, points to himself and says: ``I and the Father, we are one.''

The question how an actual, hence powerless, human being on earth can be one with the Almighty in
heaven is no question of knowledge but a personal question, hence a question of faith. If the Father
is infinite in omnipotence and freedom, then the Son is infinite in devotion and obedience: they are
different and yet one in the spirit of love, and the spirit of love is personal.

As the revealer of God's personality the Saviour stands surrounded by wonders and signs. These wonders
and signs are to make it manifest that God has not merely omnipotence and will, but a heart that feels
for the suffering and has compassion on the erring: the personally portending, the personally saving
is the core in them all. However unreasonable, indeed offensive, Christ's miracles must seem to those
who want to have everything grounded in knowledge, just so comforting and divine are they for
believers.

% --- p. 75 ---
\setcounter{footnote}{0}%
\opage{75}It might well seem as though modern criticism, which not only weakens the historical
testimonies to the credibility of the miracles but also furnishes proofs of their impossibility, must
exert a destructive, all-faith-demolishing effect on men's minds and, by constant reference to
impersonal powers, at last make the general consciousness a stranger to the personal truth. That the
movement can be dangerous enough for individuals who, weakly rooted in the personal, are easily swept
along by the current of opinion, shall not be denied; but if we look more closely, it may well turn out
that the faith-denying criticism has done positive faith in revelation a great and essential service;
it may well turn out in a time to come that criticism's thoroughgoing dissolution of the whole of
sacred history has been a purifying means, a mighty contribution to a clearer and deeper understanding
of its facts; it may well turn out that the negative spirits, far from having robbed us of the reality
of these facts, have on the contrary helped us to get our eyes opened to their proper and true
significance.

Is it then, one asks, not really true, what is assured from so many sides, that science has once and
for all proved the impossibility of miracles? We answer: No, it is really not true. A proof of the
impossibility has not been given and, as surely as human knowledge has an impassable limit, will never
be able to be given. All that science has achieved in this respect rests on a confusion of the
impossible with the improbable. Criticism has an easy time proving the improbability of a miracle, for
the miracle must after all, by its character, be the extraordinary, the surprising,
% --- p. 76 ---
\setcounter{footnote}{0}%
\opage{76}the improbable; but that the improbable is not the same as the impossible, the calculus of
probability proves even mathematically. The real matter at issue between science and religion turns,
as said, not on the reality or unreality, probability or improbability, of a historical number of
miracles, but solely and only on the possibility of the wonder. One can assume this possibility
without coming into conflict with science, for it does not lie at all within the domain of scientific
inquiry; but one cannot deny this possibility without denying God's personality and therewith the
whole of revealed religion.

To what extent what has been said holds can be seen in a survey of the evangelical miracles; with all
their externality they are all without exception to be conceived according to the law of inwardness,
from the standpoint of personality.

It is with many sincere believers in the Bible an old, ingrained prejudice that a miracle can be
experienced like any other historical fact, if only one could become an eyewitness. That is a false
presupposition. The outward, the surprising, the astonishing one can experience; but the cause? A
cunning use of secret forces of nature? A prearranged game among initiates? An illusion of the senses
or a deed of God? Here an explanation must be chosen, and the choice is personal. No miracle in the
life of Jesus is otherwise constituted. He heals one possessed, and the Pharisees say that it is with
the Devil's assistance. He raises the widow's son; had a Sadducee been present, he would have said:
``it happened conveniently, the young man was only apparently dead.'' He calls Lazarus forth from the
grave; and there are those
% --- p. 77 ---
\setcounter{footnote}{0}%
\opage{77}who think that it was a prearranged game with the friend and his sisters. It is
characteristic that some of the eyewitnesses went to the Pharisees and told them what he had done, and
that it was then resolved in the assembled council to put him to death. Such a resolution would have
been downright insane if the Pharisees had really been able to convince themselves that the miracle
with Lazarus was a historical fact: to want to put to death a man who can raise the dead is more than
defiance and unbelief, it is simply and plainly insanity.

When we read the accounts in sacred history, it seems as if each miracle stood before intuition in
such clear lines that no eyewitness could be mistaken, whether he was otherwise friend or foe. That
comes from this, that the accounts, as they lie before us, are conceived in faith and presented in
faith. Had we had the accounts of unbelief lying before us alongside, we should have seen a difference;
and yet the eyewitnesses can, as far as the historical is concerned, be assumed to be, each from his
standpoint, equally honest, equally credible.

No credible man reports anything other than what he has really seen; but the personal difference is
that he who sees a miracle with the eye of faith gets to see something other than he who sees with that
of unbelief. The account is not written down in the same instant as the miracle happens; an interval
elapses, and in the interval the outward experience passes over into the inward. It is not memory that
becomes the decisive thing; it is recollection. But the recollection of faith is other than that of
unbelief.

% --- p. 78 ---
\setcounter{footnote}{0}%
\opage{78}It is an old prejudice among believers in the Bible that the many individual miracles should
be just so many historical proofs of the truth of revelation. This misunderstanding the more recent
criticism has duly chastised. If the miracle is to stand firm as so irrefutable a proof that no one can
deny its historical reality without denying all history, then we shall, says criticism, soon get it
proved that the proof is worthless, that the whole revelation is a myth.

``What then is the miracle for, when it both steps forth and yet does not step forth, is both to be a
proof and yet not a proof?'' We answer that both determinations belong originally to the manifestation
of the personal. Did the miracle not step forth, then revelation would lack the sign of omnipotence
which says that God is personally present. Did the miracle step forth with so outwardly overwhelming a
power that there was here no possibility of a choice between faith and unbelief, it would crush human
personality: before power and the powerful all must cast themselves in the dust.

It is related that Pilate, when he learned that Jesus was a Galilean and ``belonged unto Herod's
jurisdiction,'' sent the prisoner to Herod, who himself was at Jerusalem in those same days. ``And when
Herod saw Jesus, he was exceeding glad: for he was desirous to see him of a long season, because he had
heard many things of him; and he hoped to have seen some sign done by him.'' How like a Herod it is to
want to enjoy the miracle as an entertainment and to play with the signs of power. But with the signs
of omnipotence there can be no playing; and if the captive man had lost patience, and the Almighty in
heaven likewise, in other words: if the miracle had really
% --- p. 79 ---
\setcounter{footnote}{0}%
\opage{79}come to issue in an unveiled, terrible act of power, what then would have become of Herod,
Pilate, Caiaphas and all those who had their power from crawling before power? Indeed, what would have
become of human freedom?

If the faith in the Bible that would build directly on the outward letter were right --- which it
nevertheless, by its own presupposition, cannot possibly be, since the wording of the biblical accounts
is on many points so divergent --- but suppose that the old faith in the Bible were right, then the
conditions of the life of faith would be quite other for us than they were for those eyewitnesses.
Those who immediately experienced the miracles themselves must indeed get ``faith into their hands'';
but we, who live in a world without miracles and must therefore live on old accounts, how should we get
faith into our hands? If on the other hand we hold fast to what lies in the determination of the
miracle, namely to awaken the personal life in man, then the condition for us and for those becomes
essentially the same. Just as the eyewitnesses stood face to face with the \emph{facts} with the great
question of the possibility and a personal choice between the explanation of faith and that of
unbelief, so we stand face to face with the \emph{accounts} to this very day with the same question
and the same choice.

What an enormous distance there is between a scientific grounding of a miracle's reality and a
believing understanding of the real miracle's significance can be seen by a glance at what has been
said. To introduce miraculous representations into an objective science is to hollow out its foundation
and corrupt its essence. But to bring the knower to see the possibility of the miracle
% --- p. 80 ---
\setcounter{footnote}{0}%
\opage{80}is not to encroach upon science; with this possibility as foundation to build up the cognition
of faith in scientific illumination is to reconcile the old conflict between faith and knowledge.

What above all matters in our present period of culture is this twofold thing. \emph{From the
standpoint of those who know}: a recognition that revealed religion is a spiritual power which does not
let itself be dismissed; that the religious life has a validity in itself, a personal truth which is
independent of science. The knower can have such a recognition without being a believer; but if the
recognition is lacking, then the state of culture is hollow, and all that is done from a supposedly
higher political station out of regard for the common people and the populace, to protect the
religious institutions, is wasted effort. The common people have --- in such things --- a sure instinct
and a sharp eye that does not let itself be dazzled. \emph{From the standpoint of believers}: a
recognition that the religious life in our day needs understanding, inward, essential understanding of
the revealed facts: without understanding no appropriation, without appropriation no experience.

To understanding belongs comprehending, and to comprehending thinking. Images, representations,
feelings, imaginings and opinions do not make the matter clear. Faith has its concepts just as fully as
science has its own. The concepts of religion and of faith form in their circle just as well a coherent
unity as the scientific concepts do in theirs. But that a boundless confusion in thought must arise
when the two circles of cognition are mixed, anyone can with some
% --- p. 81 ---
\setcounter{footnote}{0}%
\opage{81}reflection easily convince himself. We shall not here paint the dismay that would have to come
upon even the greatest scientific genius if he were officially to be compelled to give an objective,
incontrovertible grounding, satisfying to all, of the doctrine of God's mercy and the forgiveness of
sins. --- Let a thinker who has spent his life and his talent pondering the religious truths, who with
unconditional respect for science and scientific research has conscientiously gone through, discussed
and tested all that has been objected from the side of criticism against the Saviour of the world and
his miracles, have come to the result that the evangelists are right, that in faith there must be a
life which science does not know: has he then thereby become a believer? He can defend faith against
the attacks of knowledge, and still, personally speaking, things may look doubtful with his own faith.
There is something he has gained, and it is certainly something essential: he has an understanding in
thought. An understanding in thought is a condition, an indispensable condition for a true
appropriation, but appropriation itself rests on experience: no one comes a hair's breadth further in
faith than he comes in personal experience.

For the sake of life, that is, for the sake of experience, the Saviour of the world came into the
world. We cannot first in an objective theory assure ourselves of the validity of his coming, and then
build our faith on that. His personal life is the spring that opens itself to all. Every breakthrough of
new and higher life in history is conditioned by a human sacrifice: the Saviour is he who sacrifices
himself for all. Therefore he presents himself as man and yet with divine authority and lays decisive
weight on his
% --- p. 82 ---
\setcounter{footnote}{0}%
\opage{82}personal I: ``Come unto me, live on me, abide in me: without me ye can do nothing.'' It is the
eternal life of personality that reveals itself in him and completes its circuit: from divine origin
through abasement to exaltation.

It is remarkable to see to what degree objective knowledge has tormented itself to get the human in the
Saviour's personality to agree with the divine. Gnosticism in antiquity held to the divine and declared
his suffering and death to be semblance and illusion. Rationalism in the present holds to the human and
explains his resurrection and ascension as a myth. Only by rejecting objective knowledge, holding to the
cognition of faith and fixing our gaze on the personal does it become possible for us to see the real
connection between the abasement and the exaltation and to understand the outward by virtue of the
inward.

The outward in the history of the Passion is repellent to the natural sense, just because it is so real.
The sun of the Godhead goes down, the shadows spread, it grows so dark. The dark way of experience goes
first to Gethsemane, then to Pilate's judgment hall and then to Golgotha. The outward, surrounded by
horrors, repels, drives the beholder into himself and awakens question upon question: what is it, after
all, that is happening here? Who is the condemned man and what are the judges? Does righteousness look
like this, is this the fate of the righteous? No aid from above, no angels with swords, no saving
miracle; only reality, the stern, ice-cold reality, the laws of the world with inescapable demands for
fulfilment! The Almighty sees it from his heaven,
% --- p. 83 ---
\setcounter{footnote}{0}%
\opage{83}but does not stir; Elijah does not come down to help the helpless. The voices are silent above,
so that the words of blasphemy, of hatred, of scorn can be heard all the louder from below. Only when the
righteous one has brought his sacrifice and the thorn-crowned love has bowed its head do the portents
come: the sun is darkened, the earth quakes, the veil of the temple is rent; what does all this mean? It
means that the sacrifice is accepted, that the impotence of love has, in the name of freedom and
obedience, become one with its omnipotence; it means that the dead one shall overcome death itself, and
the personal gain mastery over the impersonal. But the outward is terrible: faith is bowed down, hope
has sorrow.

With her sorrow and her hope a Mary Magdalene, the believing soul, goes to the garden, to the place where
they laid him, and peers into the tomb: alas, he is gone. She weeps, turns around and thinks she sees a
gardener who can tell her where things stand; then she hears her name: Mary! Nothing works with greater
inwardness and more personally than to be called by one's name thus, in such a moment and by such a
voice. Here too the understanding of the outward proceeds from the inward. In the dawning morning, in the
dimness of the surroundings, shadow and reality flow over into each other; how natural that a woman's
tear-blinded eye can mistake the one her soul expected, the object of her longing and hope, for a
strange gardener. But --- Mary! At the name and the voice it dawns in her soul: she sees because she
believes, she sees in blessed certainty the fulfilment of her hope and worships the Risen One.

If the history of the Passion has something repellent because the real is so terribly real, then the
history of the resurrection
% --- p. 84 ---
\setcounter{footnote}{0}%
\opage{84}and the ascension --- for the natural sense --- has, conversely, its doubtfulness in this, that
the ideal, the super-real, does not seem to be real enough. How many times, under which and how many
forms the Risen One showed himself to his believers cannot be decided; away with bungling attempts in
objective knowledge! But it is certain that the Transfigured was never seen by the eye of unbelief; and
this is then again a proof that the outwardly real is there to awaken and confirm the personal life in
the inward. Our hope of immortality is grounded on a wonder, on the Saviour's resurrection. This wonder
must, for the sake of life and of hope, be conceived as a manifest, for believers really historical,
fact. But the believing consciousness errs when it would lean directly on outward proofs and engage
with objective knowledge.

Had it lain in the determination of the wonder that the proof of its outward reality should be
irrefutable, then the Risen One would only have needed to show himself in broad daylight in the streets
of Jerusalem. What a thronging of spectators, Jews, Greeks, Romans! there should in truth have been no
lack of eyewitnesses. Little would it have helped Sadducees and doubters to appeal to the claim that the
Crucified had only been apparently dead, that he had been taken down from the cross too early; for the
resurrection ends, after all, with the ascension. Here, then, is at last a miracle which, had it taken
place publicly and in the sight of all, could have been confirmed by so magnificent a historical
testimony that not a shadow of doubt would have remained. Imagine what an impression it would have made
on the whole population if the Risen One had appointed the day and hour
% --- p. 85 ---
\setcounter{footnote}{0}%
\opage{85}when he with his own would ascend the Mount of Olives and from there go up to heaven to sit at
the right hand of omnipotence. The multitude of men, however distracted and thoughtless it may otherwise
be, would surely not have shown itself indifferent to so magnificent a spectacle. It would have been more
than a curious crowd run together that had here come to play the role of the eyewitness. Despite the
contempt of the mighty and the scribes for the Nazarene's miracles, the real resurrection, however they
interpreted it, would already have roused their attention and put them in suspense. Like wildfire the
tidings would surely have reached Rome, Athens, the great cities of the world: the whole of civilized
humanity would have been represented at this spectacle. What a sight, to see the Transfigured on the
cloud, high above all, visible to all. Then we would have had a historical proof, a proof of fact, that
could strike every doubt to the ground.

Those who insist on the outward, irrefutable proofs must surely here become afraid of their own
representations. A proof of fact like the one described above would indeed have glorified the divine
personality by annihilating the human. When everything was fixed fast in the externality of the
phenomenon of power, what then would become of the free choice between faith and unbelief? What would
become of personal development? It would have been annihilated; the whole would have ended with a Day of
Judgment.

The Transfigured stands on the cloud; let his countenance be ever so mild, the fact speaks: ``Ye would
not believe, now ye have certainty, outward certainty, and the outward certainty judges and condemns you
all. Ye would
% --- p. 86 ---
\setcounter{footnote}{0}%
\opage{86}not have freedom in obedience; now freedom departs, it goes up to heaven, whence it came; but
ye shall bear the yoke of obedience, ye shall slave in the dust. Ye scorned love when it was seen in its
impotence; ye mocked its self-sacrifice, its humble patience, its rich promises, its anxious care, its
heartfelt compassion, even when it was seen in tears. Now ye shall see it clad in its omnipotence, the
Majesty as he is, not in the form of the servant, no, in that of the ruler: he crushes every powerless
defiance and asks no more after the will of men.'' That would be a proof of fact, but a terrible one.

If we consider this, then we can understand why the Risen One came so quietly, as in secret: to gather
with his little circle, to walk with the two to Emmaus, when the day was far spent and it was toward
evening, to greet the eleven when the doors were shut. Suddenly he stands among them, as a vision, as a
spirit, and yet as he really is, the Crucified himself with the open side and the print of the nails in
his hands. It is the way of inwardness, the quiet way; not that reality should be doubtful, but that the
testimony may come from without and yet from within, from the depth of the soul, from the ground of the
mind, that the manifestation may be not merely for sight and sense, but for inner experience unto
personal growth, for true appropriation in faith and hope.

% --- p. 87 ---
\setcounter{footnote}{0}%
\lecture{Sixth Lecture.}\opage{87}Every age has its prejudices, our own no less than all those that went before. The prejudice
of the Middle Ages was that everything concerning faith and hope had to be confirmed from without by an authoritative
decree; its faith in the divine authority of the Church and the Pope was a superstition. The prejudice of the present
is that a religious conviction is without value when it has not received its baptism in an objective knowledge; this
faith in the ability and the exclusive right of science to decide what may pass and count as religious truth is the
superstition of our time. Both errors, the modern as well as the medieval, show themselves in their true light
especially when the talk is of the completion of life, of ``the last things,'' of the state after death.

The Middle Ages were almost more topographically oriented in the world beyond than in this world. One had all the
circles before one: Paradise, the circle of the Fathers or Abraham's bosom, the circle of the children, Purgatory and
Hell; one knew exactly where it all lay. Protestantism struck a line through the circle of Purgatory, and kept on the
whole to the Bible. Here, then, one thought to find reliable reports from
% --- p. 88 ---
\setcounter{footnote}{0}%
\opage{88}the other world. Christ, who had come down from heaven, where he had seen everything with his own eyes, must
surely know how matters stood; the Apostles, who had every single jot they wrote down inspired by the Holy Ghost, could
then not be mistaken either. So one collected passages of Scripture; but the more eagerly one strove to work the
scattered hints together into an intelligible whole, the greater the difficulties became.

The parable of Lazarus and the rich man, Christ's words to the thief on the cross and similar utterances seem, after
all, to put it beyond all doubt that there is no intermediate state: Lazarus comes at once into Abraham's bosom, the
rich man into Hell, the thief into Paradise. But in opposition to this stands the doctrine that the dead shall not rise
from their graves until the last day. How then does it stand with the soul in the interval that lies between death
and the resurrection? Some among the ancients thought that the soul died with the body, and that both, soul and body
alike, were awakened to life again by a miracle; others thought that the souls slept until the Day of Judgment, and
still others, who rejected the sleep of souls, found the intermediate state designated by laying weight on the sojourn
of the spirits in Hades. --- All of it bungling in objective knowledge.

The endeavor to get the doctrine fixed dogmatically in objective forms stirred up reflection and brought on
criticism. It was the half-rationalism of the last century that broke the ice. With incredible racking of brains one
exerted oneself and wanted to make it quite plain to sound human understanding how it really must go on the Day of
Judgment. The event was divided juridically into three sections: 1) the examination,
% --- p. 89 ---
\setcounter{footnote}{0}%
\opage{89}2) the passing and the publication of the sentence, and 3) its execution. Then it dawned on the
understanding that the whole could not possibly be finished in one day. ``Indisputably,'' says an older dogmatician,
``the judgment of the world must last long, which is all the more reasonable as a clear presentation of this important
matter will be the most instructive spectacle for every thinking being and spread the greatest light over the darkest
mysteries. Unfortunately the external circumstances cannot be more closely determined.'' According to a probable
estimate it has been found that for a number of about 300 million human beings the Day of Judgment must last at least
100,000 years.

Another difficulty was: to find out how so many human beings could at once come to see the Judge and hear his words.
Some assumed that the risen bodies were free of space; but it was easy to prove that they could not possibly be bodies
without occupying a space. Even if one would hold that, if they really could squeeze themselves together and make
themselves small, a thousand could sit on the head of a pin, still there would after all have to be room. As late as
the nineties an anonymous author worked out that if human beings, as is only fair, were to stand in court, an area of
the size of the Kingdom of Bohemia would be sufficient without squeezing.

Those who cannot see here that faith and knowledge must go each its own way should be loath to take it upon themselves
to be guides. The dogmatics of the Church has, partly historically, partly speculatively, gone into partnership with
objective knowledge: it has not agreed with it. Objective knowledge wants, as lies in its nature, to know how matters
stand
% --- p. 90 ---
\setcounter{footnote}{0}%
\opage{90}with the objective; it relates itself --- which, from a half-speculative dogmatic standpoint, has indeed
been boasted of in recent times as well --- quite rightly to the supernatural phenomena as the natural scientist
relates himself to the natural facts. But when objective knowledge has once become master in the house, what then? By
subjecting the supernatural phenomena to a criticism it discovers that the whole is absurd, incoherent, improbable and
contradictory; with this discovery it grows angry, smashes the ecclesiastical porcelain of dogmatics to pieces and
goes its way.

Otherwise with the cognition that springs from the principle of faith; it does not concern itself with investigations
of the objective, it applies another standard. The first question is: what significance have the utterances of
Christ and the Apostles about ``the last things'' for the personality itself and the personal experience? It then
becomes at once evident that these utterances do not want to satisfy a human curiosity, but to guide us to understand
what is decisive in an eternal sense. The revealed word wants in holy earnest to awaken our fear and our hope; it wants
with divine authority to impress upon us representations of the great crises that life has to pass through before it
reaches completion. From this point of view we can understand that the different utterances, although they all spring
from the same fundamental thought, can proceed from mutually different motives. The motive of the presentation
changes each time the same truth is to be illuminated from a new side.

I. In the parable of the rich man and Lazarus, as also in the promise to the thief on the cross, it is
% --- p. 91 ---
\setcounter{footnote}{0}%
\opage{91}expressly stated that between the life here and the life yonder there is an essential connection, an inner
personal coherence: the coming state is prepared and conditioned by the direction taken in the present. The poor man
comes into Abraham's bosom as one who has suffered in faith and patience; the rich man with his egoism of enjoyment and
his mercilessness gets his place where mercilessness has its home. The thief, who in the hour of torment receives the
vision of faith and discovers that his fellow-crucified has a kingdom in which he is Lord, comes to the place of peace
and blessedness, and the place is indeed quite near him.

Another motive lies in the presentation of the resurrection. Here the individuals go together with the race, and the
separation between the good and the wicked is accomplished. Just as there is one death for all, so there is also one
resurrection and one judgment for all. It is an earthly eye that is to be opened to see the supramundane. If we now
take our stand at an earthly standpoint and ask: where are our dead? then the natural answer is: they rest in the bosom
of the earth, they sleep in the graves. It is a real view of the dwelling-place of the dead, an impression of the
perishableness of life; the power of death is imposing. To this there must then answer an equally grand, imposing
impression of the power of the vanquisher of death, of the power of him to whom it is given to have life in himself.
Supramundane powers stir, the trumpet of judgment sounds, the voice of the Son of Man is heard: the sleepers awake, the
graves open, the dead rise, the Day of Judgment begins. Although in all this the real is what is aimed at, the
presentation nevertheless keeps to the form of the symbol; it lies in the nature of the matter, and is immediately
discerned from it, that the images alternate.
% --- p. 92 ---
\setcounter{footnote}{0}%
\opage{92}First the race is judged, in that the generations judge one another. ``The men of Nineveh shall rise in
judgment with this generation, and shall condemn it: because they repented at the preaching of Jonas; and, behold, a
greater than Jonas is here. The queen of the south shall rise up in the judgment with this generation, and shall
condemn it: for she came from the uttermost parts of the earth to hear the wisdom of Solomon; and, behold, a greater
than Solomon is here.'' That the standard is taken from generation against generation is a testimony that the judgment
comes from within, that it is not arbitrary, but personal. Then another image: around the throne of Christ sit twelve
judges, the Apostles on twelve seats; they judge the twelve tribes of Israel. Finally the depiction is completed: the
Son of Man is seen with all the holy angels; he sits upon the throne of his glory: ``before him shall be gathered all
nations: and he shall separate them one from another, as a shepherd divideth his sheep from the goats.''

The fundamental concepts that bear the possibility of the whole supramundane content all have the stamp of full
reality. The two worlds, this world and the world beyond, both have space, time, matter and forces, the one in
perishable, the other in imperishable form: nothing inward without something outward, no souls without bodies, no
human personality without real, visible individuality. Feasts are held in the world beyond: they come from east and
west to sit at table with Abraham, Isaac and Jacob; Christ drinks with His own --- the new wine in the Father's
kingdom. These are forms of imagery, full of the content of truth; but it is no description of an external, fixed
reality, no map of the foreign country, no
% --- p. 93 ---
\setcounter{footnote}{0}%
\opage{93}photograph of its regions, its buildings, its household furniture, its inhabitants.

What is it, now, that distinguishes the great images of revelation from all myths, all forms of apocalyptic poetry
and fantastic symbolism? It is, in a single word, the personal. The fundamental thought, the kernel in it all, is the
completion of the personal life. Everything impersonal is means, material, apparatus. It is the personality itself,
the divine and the human, the \emph{divine-human}, about which everything turns. This axis of rotation is firmer than
the astronomical axis of the world, this celestial point is more unchangeable than the pole star.

It is a misunderstanding in which only thoughtless superficiality combined with a habitual propensity for systematic
befogging of the brain can take pleasure, an altogether absurd misunderstanding: that the referring of all miracles and
supernatural phenomena should coincide with a cunning tendency, a tacit endeavor to get the reality of the
supernatural, i.e.\ of the higher natural, explained away, or at any rate transformed into an indeterminate idea
floating in the air. On the contrary, it is surely in an endeavor to assert the character of this reality as an
indissoluble spiritual reality, unassailable by science, that the personal point of view has been chosen. Why do we
reject the bungling with objective knowledge? In order to secure the standard of the personality, and with it of
experience, for faith and hope.

Faith is a life; as the life is, so is the hope, and as the hope is, so is the life. If there is an essential
difference between the life of faith and the selfishly natural life, then there must surely also be a corresponding
% --- p. 94 ---
\setcounter{footnote}{0}%
\opage{94}difference between the hope reborn with faith and all the forms of natural hope that touch the life of faith
only from without. It is, when we look a little more closely, not difficult to separate natural hope from the hope of
faith. The conception of the relation between the present and the coming world affords a sure test. Natural hope
resembles the hope of faith in that it always wants something higher, more beautiful, more perfect; so long as gradual
change suffices, so long as facts and representations borrowed from familiar surroundings can support the
probability, natural hope goes along, even if the bridge that leads over the abysses is narrow and long. But if the
railing of probability falls, if all that is familiar is to be given up, and the transition is designated by a
transformation, then it grows dark before the eyes of hope: it either becomes wildly fantastic, or it grows faint and
dies. Only the reborn hope of faith endures the transition of transformation. It is only this hope for whose
fulfilment the revealed word will vouch; every other hope attaching itself to religion can be disappointed.

It is remarkable to see how easily human consciousness can mix together what is absolutely heterogeneous. There is in
many souls a longing, a hope, that is built on faith; they hope for a resurrection, a reunion yonder, a
transfiguration, an eternal blessedness. Faith would so gladly hold fast to all this; but is it not strange? Hope,
which is otherwise so willing to hasten ahead and strew gleams of light on the way, here becomes involuntarily
frightened and dares not go along. What then is in the way? ``We must, after all,'' it is said, ``necessarily believe
in an eternal life and a reunion yonder; if we did not
% --- p. 95 ---
\setcounter{footnote}{0}%
\opage{95}hold fast to this faith, the present life would be worse than death. But it is nevertheless so difficult for
us to believe, for hope is so faint, it will not shine. If we could only form a tolerably reasonable representation
of what awaits us beyond, and of how the real transition from death to life is actually to come about, then hope would
strengthen our faith.''

Instead of referring to the transformation and demanding the rebirth of hope, there now present themselves the
exceedingly agreeable, most obliging guides, willing to support natural hope with a little objective knowledge and to
make the transition comprehensible, highly reasonable. Hence these half-naturalistic, half-believing flirtations with
``the day after death,'' partly with hydrogen and ether, partly with Bible passages, presentiments, assurances and
modest, yet unpresuming opinions. One forgets that the transition is a transformation, that only the reborn hope can
appropriate the transformation and go together with faith.

When the Saviour, on the last evening, at the farewell, comforted the sorrowing disciples, his words were: ``Let not
your heart be troubled: believe in God and believe in me. In my Father's house are many mansions\dots{}. I go to
prepare a place for you.'' The sentimental half-faith with its natural hope finds it highly edifying and therefore
quite acceptable that with these words Christ pointed up toward the visible evening sky, toward Sirius, the Pleiades,
Orion, and, who knows, perhaps even prophetically gave a review of ``the Manifestation of Grace in Nature.'' The
insipidity of such flirtations lies in the feeling that there is no earnest about either nature or grace. No one has
in stronger and more authoritative words
% --- p. 96 ---
\setcounter{footnote}{0}%
\opage{96}declared that the whole visible universe is subject to perishableness than the Saviour himself, when he % print: har (misprint)
says: ``Heaven and earth shall pass away: but my words shall not pass away.'' So long as Christ's words have any
meaning, the globes in astronomical space cannot be assumed to be mansions for the blessed.

Entrance into ``the eternal mansions'' is conditioned by a transformation. It is the transformation, not a distance in
space, that sets the boundary between the present and the coming world. Let us consider what a transformation means.
What change is in the form of finitude, transformation is in that of infinity. Grief, sickness, old age can change a
human being even to such a degree that those who have not seen him in the meantime find him quite unrecognizable; and
yet the whole is only a change, inasmuch as it falls within the bounds of finitude. Conversely, a transformation can
take place in a human being without the person thereby becoming unrecognizable. A young girl, a youth, a mature man
can experience something that so completely refashions the inner being that we must say: it is the same human being,
and yet not the same. All merely partial and gradual transitions, however far they go, are changes, but the
qualitative transitions are transformations. That a human being little by little lays aside his faults is an
improvement, and an improvement is a change; but the transition from unbelief to faith is more than a change, it is a
transformation; rebirth is a transformation; the transition from life to death and from death to life is a
transformation; Paul's conversion is more than a change, it is a transformation.

% --- p. 97 ---
\setcounter{footnote}{0}%
\opage{97}Only when the difference between change and transformation has become clear to us can we, on the terms of
faith, form a concept, namely a religious concept, of the relation between a perishable and an imperishable world.
Light and shadow, laws and limits, stamped forms and full reality must surely be assumed to belong just as fully to
the imperishable nature as to the perishable. And yet here there is a thoroughgoing difference. There is in the
shadow, in the limit, in the essence of difference and opposition, a two-sidedness that must not be overlooked. It is
not the shadows, but their disproportion to the light, that testifies to the perishable: the true proportion between
shadows and light gives us, after all, precisely the richness of colors. It is not the limits in and for themselves,
but the disproportion between the limits and that which is limited, that characterizes the perishable. In the form of
imperishability the limits are not barriers, checks, hindrances, but determinations of essence for the individual.
Even pain, suffering, resistance, sorrow, fear, danger receive the stamp of perishableness solely by stepping forth as
independent powers that subdue and threaten the living being. The imperishable nature does not simply obliterate these
negative powers, but it transforms them, reduces them to overcome hindrances, serving determinations for the
positive, full living. Only in the perishable world do the dissonances have a shrieking independence; in the
imperishable they sound along, but this sounding-along is a resolution into harmonies: the harmonies themselves have
a continued resolution of the dissonances as their foundation. A true cognition of what transformation is presupposes
experience; the reborn
% --- p. 98 ---
\setcounter{footnote}{0}%
\opage{98}hope is built on faith in inner experience. It is to this and only this hope that the Revealer of the
universal love has given his word, and set his life in pledge, that it shall not be disappointed.

If we consider the personal relations of the individuals among themselves, relations which have a psychical, spiritual
stamp of eternity that raises them above the perishable, then the difference between the natural and the reborn hope
will at once become visible. There are bonds of love that bind parents to children, the husband to the wife, the lover
to the beloved. Death tears these bonds asunder; what pain! It is indeed as if it tore the personal up by the roots,
scattered it to all the winds, gave it up as prey to impersonal elements. Yet there is, after all, a reunion! Faith
calls on hope, but hope despairs: what became of the form, where is the land of reunion, where are the possibilities?
Natural hope must die in order to be born anew, and of what shall it be born if not of faith in him who hath eternal
life in human form and shape. His resurrection and transfiguration is the only guarantee we have that individuality is
not dissolved when it is transformed and transfigured. ``Ye now,'' he said to His own, when death stood just waiting
for him, ``ye now have sorrow: but I will see you again, and your heart shall rejoice, and no one shall take your joy
from you.'' When he saw his own again, and they saw him, he was transformed. But was he therefore less real, less
individual, less personal?

The hope that is born of the power of his resurrection in personal experience forms the foundation for every
particular, individual hope of a reunion. None of the
% --- p. 99 ---
\setcounter{footnote}{0}%
\opage{99}personal features is erased, nothing of the form on which the eye of love has rested is lost; only one thing
is lost, irrevocably lost, and that is the perishable.

And then, finally, the goal of completion, blessedness; here the same thing repeats itself: faith wants to hold fast
that there is a blessedness yonder, but natural hope loses heart. What is blessedness? A word that designates
transformation. If it could begin straightforwardly with what we know, with the happy, and the happy rise to the
blissful, and the summit of bliss reach right up to blessedness: then natural hope would spread its wings with
delight. Natural hope understands so well how to give the happy, the perishable, a little spicy admixture of the
imperishable, the more than happy, the blessed, and to let both elements go together into one in bliss. That two who
have long struggled with dangers and adversity and, what is worse, with mutual misunderstanding, must in the moment of
union and understanding become not merely happy, that is too little, no, blissful --- that natural hope can understand
excellently; it breathes deeply, and the heart beats at the representation of the bliss of the two. But to their
blessedness, if they succumbed in the struggle and died on the way to the goal of bliss, the disappointed hope would
offer only a tear.

Natural hope does not grasp eternal blessedness; for eternal blessedness does not enter into alliance with happiness,
any more than imperishability in the world beyond enters into alliance with perishableness. Happiness must go, for
perishableness must go; now it grows dark before the eyes of hope. Blessedness, pure
% --- p. 100 ---
\setcounter{footnote}{0}%
\opage{100}sheer blessedness, what kind of state is that, after all? Light in light and white in white, all images
grow so pale: the white robes, the palms, even the harps and the hymns of praise must surely, because of the
everlasting repetition, lose themselves in monotony. If natural hope is not to suck the blood of faith and end in
despondency, it must die; but it dies only to be born anew.

The hope of blessedness, the reborn hope, is not faint; it is in a growing tension. It receives, for personal
reasons, an opposite that will surely keep it awake; this opposite is fear. What, then, has he who hopes for
blessedness to fear? That which is opposite to blessedness. It holds in the imperishable as in the perishable that the
opposite terms must correspond to each other. What is opposite to happiness? Unhappiness. What is opposite to bliss?
Misery. And what, then, is opposite to blessedness? Unblessedness, perdition. From this it is seen that fear, in a
special sense, must grow with hope. The fainter the hope, and the paler the image of blessedness, the less too is the
fear of unblessedness. Conversely, the richer and stronger the image of blessedness becomes, the more inward the desire
of longing that unites itself with hope, the greater too is the involuntary dread of the possibility of
unblessedness. Fear grows at the thought that on no point is it easier to deceive oneself than precisely on this.

There is a disquiet of mind that can torment the soul with uncertainty and inflict on it all the agonies and torments
that made the medieval asceticism and the pietistic darkening of the spirit only too easily explicable. Against these
excesses the Lutheran thought stands as a sheltering wall: ``Thou shalt believe, not in thine own, but
% --- p. 101 ---
\setcounter{footnote}{0}%
\opage{101}in the Reconciler's righteousness; in faith thou hast the certainty of blessedness, a certainty that cannot
be shaken!'' Are we then done with this? By no means! Just as faith wants to fix itself as something that is,
something settled once for all, deception steals in from another side. ``We are believers, and in faith hoping; with
us there is no danger. We know no fear, we know only joy; we are like children, glad of the present, gladder still of
what is to come, glad on earth, gladder still in heaven, glad of happiness, gladder still of blessedness. The darkness
is gone, the prospect clear, the rod is broken, all is a game: so gladly do we play before our Lord.''

Truly, if the Spirit is not so angered by such an offence against the revealed word that it turns away and lets us
remain in our pleasant imaginings, with bold courage to make a noise in society, hoist the victory flag of
thoughtlessness and go adventuring with those who work havoc, then the holy Spirit of truth must indeed turn against
us in flaming wrath. This wrath will then be the testimony to its incomprehensible mercy; for it will ask: ``what is
your hope, where is your fear? I --- the Holy Ghost --- know only one hope, the hope that dies in temporal
disappointments, but rises with the Risen One, the hope that does not run away from the inner dangers, but overcomes
all fear, the reborn hope. What is your freedom, where is your obedience? I know only one freedom; it received its
baptism in self-denial, its consecration in obedience. What is your joy, where is your sorrow? I know a joy
unspeakably great; the starry heavens cannot contain it, and yet it can dwell in a human being's breast; but it is
not a babbling
% --- p. 102 ---
\setcounter{footnote}{0}%
\opage{102}children's joy, it is the joy that can lift sorrow, drain the cup of pains and say to the cross: Welcome!
Where is the struggle of life that ye withstood? where is your great will, the high mind, the clear deed, the sign of
eternity by which the heavenly Father wills that his children shall be known in the world? What is your faith, your
hope? show me your transformation!''

But, objects the reflecting understanding, if there really is an imperishable world in store, what then is the use of
these scruples about transformation, so long as we are in this perishable state; the transformation surely takes place
of itself when the transition to the imperishable occurs. What does this dread of unblessedness and perdition mean? In
an imperishable life there can surely be no unblessedness. Here it is forgotten that the kingdom of blessedness is a
kingdom of personality and freedom, that the peace and the rest lie in the harmony of the active wills with the
all-fatherly fundamental will. If there are now wills that do not merge into this harmony, souls so dark, so faint and
dull, that the great impressions have a stupefying effect; souls so selfishly hardened that they are filled with
hatred of what in the name of love is noble, exalted and pure: then surely all the rays from the sun of eternity will
burn in these souls as a consuming fire: blessedness will for them all transform itself into unblessedness.

There are, then, in that kingdom of imperishability possibilities for an infinite series of crises of will. From
below, in the dark deep, in the abysses of eternity, a defiant, negating will that stirs up everything selfish,
everything egoistic, and incites to desperate struggle amid convulsive contradictions. From above, the image of the
Redeemer, the manifestation of the Saviour in millions and millions of
% --- p. 103 ---
\setcounter{footnote}{0}%
\opage{103}reflections, and always the same voice, the same calling voice: ``come unto me, live in me, abide in me;
without me ye can do nothing.'' --- Such a peace is more than an idle rest.

Are these crises, these tensions, without end? No! Only, they end not in the manner of finitude, but in the manner of
infinity. That something ends in finitude means: it ceases, and then it is over, perhaps forgotten; one passes on to
something else. But that which ends in eternity --- its end is beginning, and its beginning is end. Here come the great
images of life, the images of hope and of fear. Everything gathers as in one vision, in one now, so that the
supramundane can become visible from an earthly standpoint; so that we must appropriate the Saviour's words: ``the
hour is coming, in the which all that are in the graves shall hear the voice of the Son of man, and shall come forth;
they that have done good, unto the resurrection of life; and they that have done evil, unto the resurrection of
judgment;'' so that we must learn to hope in the measure that we learn to fear, so that we must in truth learn to
live.

What has been communicated here is only a silhouette; and even if what has been fleetingly hinted at received a
development in ever so rich, ever so clear thoughts, received all the luxuriance of unfolding, the fertility of
expression, the persuasive power of the word in the perfection of presentation, which a higher gift can give: it would
still be only a silhouette! That which is to fill out the silhouette, that which is to raise it to reality, breathe
into it the breath of life and give it light and colors, everyone must add from his own. Life wants to be lived, hope
wants to hope; only he who lives in hope obtains the certainty of experience, the fullness of impression, the power of
conviction, and bears witness to life and hope.

\end{document}
